% AI-produced research contribution; see report.md for authorship and audit trail.
\label{partD}

\subsection{Conventions and the first obstruction over a base}
Unless another base is explicitly specified, the computations below concern a
simple adjoint algebraic group $G$ over an algebraically closed field $k$ of
characteristic zero.  Write $r=\operatorname{rank}G$, $W=N_G(T)/T$,
$\theta=\sum_{i=1}^r m_i\alpha_i$ for the highest root, $\alpha_0=-\theta$,
$m_0=1$, and $h=1+\sum_i m_i$.  The symbol $Z$ means the centre of the
simply connected covering group, and $z=|Z|$; it does not mean the centre
of the adjoint group $G$.  A splitting always means a group homomorphism
section of the specified quotient homomorphism.  An abstract bijection of
sets is insufficient.

\begin{saresult}{D.001}{Refuted}{The Lie-algebra construction is not a simultaneous resolution over every base}
The assertion that the morphism
\[
 p:\widetilde{\mathfrak g}=\{(x,B):x\in\operatorname{Lie}B\}
 \longrightarrow \mathfrak g\times_{\mathfrak t/W}\mathfrak t
\]
is birational and an isomorphism over regular semisimple elements for every
adjoint semisimple group over an arbitrary normal base is false.  It fails
for $G=\operatorname{PGL}_{2,\mathbb F_2}$.
\end{saresult}
\begin{proof}
Put $k=\overline{\mathbb F}_2$.  The Lie algebra is
$\mathfrak g=\operatorname{Mat}_2(k)/kI$; every class has a unique
representative $X=\left(\begin{smallmatrix}a&b\\c&0\end{smallmatrix}\right)$.
The diagonal torus has Lie algebra $\mathfrak t=k[\operatorname{diag}(1,0)]$.
Its nontrivial Weyl element sends $\operatorname{diag}(a,0)$ to
$\operatorname{diag}(0,a)$, which represents the same class because the
difference is $aI$.  Thus $W$ acts trivially on $\mathfrak t$, and the map
$\mathfrak t\to\mathfrak t/W$ is the identity.  The invariant trace
defines $\chi:\mathfrak g\to\mathbb A^1_k$, $[X]\mapsto a$ and extends
the restriction to $\mathfrak t$.  Consequently the displayed target is
just $\mathfrak g$.

For $a\ne0$ the polynomial $u^2+au+bc$ is separable.  Its two eigenspaces
are exactly the two invariant lines in $k^2$, hence exactly the two Borel
Lie algebras containing $[X]$.  This conclusion is unchanged upon adding a
scalar matrix.  Thus $p$ is generically of degree two, not birational.
At $\operatorname{diag}(1,0)$ the two lines are already defined over
$\mathbb F_2$.  The complete three-line enumeration is in
\texttt{code/small\_counterexamples.py}, with its output in
\texttt{logs/small\_counterexamples.txt}.  The generic assertion follows
from the separability calculation and does not depend on a point count.
This is precisely the characteristic-two issue explicitly left visible
in 123/007; it must not be removed from a modern statement.
\end{proof}

\subsection{Coxeter normalizers in every simple Weyl group}
\begin{saresult}{D.002}{Proved}{Splitting of the Coxeter normalizer}
Let $W$ be an irreducible crystallographic Weyl group and $c\in W$ a
Coxeter element of order $h$.  Conjugation gives an exact sequence of finite
groups
\[
 1\longrightarrow\langle c\rangle\longrightarrow
 A:=N_W(\langle c\rangle)\xrightarrow{\nu}
 U_h:=(\mathbb Z/h\mathbb Z)^\times\longrightarrow1,
 \qquad wcw^{-1}=c^{\nu(w)}.
\]
It splits in types $A_r,B_r,C_r,D_r,G_2,E_7$ and does not split in
types $F_4,E_6,E_8$.  Use $r\ge1$ for $A_r$, $r\ge2$ for $B_r,C_r$,
and $r\ge4$ for $D_r$; the usual lower-rank identifications cover the
remaining names.
\end{saresult}
\begin{proof}
For $A_{m-1}$ identify the letters of $S_m$ with $\mathbb Z/m\mathbb Z$
and let $c$ be translation by one.  An element conjugating $c$ to $c^a$
must have the form $j\mapsto aj+b$.  Conversely every such map, with
$a\in U_m$, is a permutation and has that conjugation property.  The
maps $j\mapsto aj$ constitute a section.  This proves both exactness and
splitting, including $m=2$.

For $B_r$ and $C_r$, realize $W$ as the signed permutations.  A negative
$r$-cycle acts regularly on the $2r$ signed letters, identified with
$\mathbb Z/2r\mathbb Z$, and change of sign is translation by $r$.
Every unit $a$ modulo $2r$ is odd, so $j\mapsto aj+b$ commutes with
change of sign.  These, and only these, are the normalizer elements.
Again multiplication by $a$ gives a section, and the centralizer is
the cyclic translation group.

For $D_r$ put $m=r-1$.  Choose $c$ to be a negative $m$-cycle on the
first $m$ coordinates and change of sign on the last coordinate.  This
has even total sign parity and order $2m$.  Its two signed-letter orbits
have lengths $2m$ and $2$, which are different since $r\ge4$.  A
normalizer element must preserve these two orbits.  On the first it is
affine $j\mapsto aj+b$; its sign on the last coordinate is then uniquely
determined by even sign parity.  This proves that the centralizer is
$\langle c\rangle$ and that every $a\in U_{2m}$ occurs.  To split, use
$j\mapsto aj$ on the first orbit and the unique parity correction on the
last coordinate.  Sign parity is a multiplicative character on the
signed permutation group, so these corrections respect products; the
result is a homomorphic section.

For the exceptional groups the following is a complete finite proof,
not a sampling of representatives.  The attached program
\texttt{coxeter\_normalizers.py} uses the displayed integral Cartan
matrices stored in its output.  It generates the complete root system
$R$ from the simple roots under the simple reflections $s_i$ and forms
$c=s_1\cdots s_r$.  It finds a root $v$ whose $c$-orbit spans
$\mathbb Q^r$, and records a basis
$B=(c^{j_1}v,\ldots,c^{j_r}v)$.  For each $a\in U_h$ and \emph{each}
$u\in R$, there is at most one possible intertwiner $w$ with
$wv=u$ and $wc=c^aw$: its matrix must be
\[
 (c^{aj_1}u,\ldots,c^{aj_r}u)B^{-1}.
\]
Thus this finite list exhausts the normalizer.  The program retains a
candidate exactly when it is integral, preserves the Gram matrix and
the entire root set, has $wc=c^aw$, and belongs to $W$.  The last test
reduces $w\rho$ to the fundamental chamber by simple reflections and
requires the resulting matrix to be the identity; the entire reduction
word is retained.  A regular dominant vector has trivial stabilizer in
$W$, so this is an exact membership test.  The recorded centralizer
equals the list $1,c,\ldots,c^{h-1}$, and every unit has exactly $h$
lifts.  This proves all three exactness assertions.

The independent verifier \texttt{verify\_certificates.py} reconstructs
the roots and all intertwiners with Python rational arithmetic, without
NumPy or SymPy, checks every reflection word, and uses a different
chamber-reduction order.  Its complete output is retained.  The relevant
results are
\[
\begin{array}{c|r|r|r|l}
\text{type}&|R|&h&|A|&\text{lift certificate}\\\hline
G_2&12&6&12&a=5:\;w^2=1\text{ for all six lifts}\\
F_4&48&12&48&a=5:\;w^2=c^3\text{ or }c^9\ (6\text{ each})\\
E_6&72&12&48&a=5:\;w^2=c^3\text{ or }c^9\ (6\text{ each})\\
E_7&126&18&108&a=5:\;w^6=1\text{ for all eighteen lifts}\\
E_8&240&30&240&a=7:\;w^4=c^5,c^{15},c^{25}\ (10\text{ each}).
\end{array}
\]
The units $5\bmod6$ and $5\bmod18$ generate their unit groups of
orders two and six, respectively, so the indicated lifts give sections.
An involution $5\bmod12$ cannot have a noninvolutory image under a
section, nor can the element $7\bmod30$ of order four have an image
whose fourth power is nontrivial.  Hence the other three sequences do
not split.  Matrices of the sections and of \emph{every} normalizer
element are in \texttt{logs/normalizer\_TYPE.json}; their presence is a
reproducible certificate of this finite proof.
\end{proof}

\subsection{The two-torus extension and the extra symplectic obstruction}
Choose a pinning, a primitive $h$th root $\zeta\in k$, and put
$t=\rho^\vee(\zeta)\in T$.  Choose a cyclic vector
$x=\sum_{i=0}^r x_i$ with $0\ne x_i\in\mathfrak g_{\alpha_i}$, and
let $T'=C_G(x)$, a maximal torus.  The Coxeter-apposition theorem says
that $t$ normalizes $T'$ and acts on it by a Coxeter element $c$.
Here and below ``in apposition'' means a pair conjugate to such a pair.
These characteristic-zero input facts are Kostant's cyclic-element
theorem and principal-element theorem \cite[\S\S6.2--6.7,
Theorem 8.6]{D-Kostant1959}.  No assertion about arbitrary residual
characteristic is hidden in this definition.

\begin{saresult}{D.003}{Proved under stated hypotheses}{The fibre product description of the common normalizer}
With the characteristic-zero apposition data just specified, define
\[
 B_t=\{w\in W(T):w(t)=t^a\text{ for some }a\in U_h\},
 \qquad\Phi=N_G(T)\cap N_G(T').
\]
There are exact sequences
\[
1\to Z\to B_t\to U_h\to1,\qquad
1\to\mu_h\times Z\to\Phi\to U_h\to1,
\]
and the two Weyl projections identify $\Phi$ with
$A\times_{U_h}B_t$.  Consequently the last sequence splits precisely
when both factor sequences split.  Type by type it splits in
$A_r,B_r,D_r,G_2,E_7$, and in $C_r$ if and only if $4\nmid r$.
It does not split in $F_4,E_6,E_8$ or $C_{4m}$.
\end{saresult}
\begin{proof}
First $T\cap T'=1$.  Indeed an element of $T$ centralizing $x$ acts
trivially on each of its nonzero, distinct root components, hence all
simple-root characters have value one.  These characters form a basis
of the character lattice of the adjoint torus.

The eigenspace of $c$ on $\operatorname{Lie}T'$ with eigenvalue $\zeta$
is the cyclic line $kx$.  An element of $T$ normalizing that line must
satisfy $\alpha_i(s)=\lambda$ for every $i\ge1$ and
$(-\theta)(s)=\lambda$.  Thus $\lambda^h=1$, and the element is
uniquely $\rho^\vee(\lambda)$.  Therefore
$T\cap N_G(T')=\langle t\rangle$: a Weyl-centralizing element must
act as a power of $c$ by D.002, and after subtracting that power it lies
in $T\cap T'$.  Conversely the displayed cyclic group normalizes $T'$.
In particular every element of $\Phi$ normalizes $\langle t\rangle$.
Its two Weyl projections have the same exponent $a$.

Let $w'\in A$ have exponent $a$, and lift it to $n'\in N_G(T')$.
Then $n'tn'^{-1}=s't^a$ for some $s'\in T'$.  Since $1-c^a$ has no
zero eigenvalue on the rational character space, the homomorphism
$T'\to T'$, $u\mapsto u c^a(u)^{-1}$, is an isogeny and is
surjective over algebraically closed $k$.  Multiplying $n'$ by a
suitable $u$ makes $n'tn'^{-1}=t^a$.  The identity component of the
centralizer of $t$ is $T$: no root has height divisible by $h$.
The adjusted $n'$ therefore also normalizes $T$.  This proves
surjectivity $\Phi\to A$; its kernel is $T'^c$.

We spell out the size and type of this kernel, to avoid replacing a
group statement by a numerical one.  The point $\rho^\vee/h$ is in the
interior of the fundamental affine alcove and has all affine simple
root coordinates equal to $1/h$.  Its stabilizer in the extended
affine Weyl group $P^\vee\rtimes W$ consists precisely of the alcove
stabilizer $\Omega=P^\vee/Q^\vee$: an element fixing an interior point
preserves its unique alcove, and every element of $\Omega$ permutes
the affine simple roots, hence fixes this point.  Its image in $W$ is
$\operatorname{Stab}_W(t)$; projection is injective on a finite point
stabilizer.  This is an abelian group isomorphic to $P^\vee/Q^\vee$,
whose abstract finite group type equals that of $Z$.
The kernel $T'^c$ injects into this stabilizer because $T\cap T'=1$.
Its order is $\det(1-c)=z$ (in classical types this follows from the
cycles above, and in exceptional types from the recorded integral
matrices).  Hence it is the whole stabilizer.  We obtain
$|\Phi|=hz\varphi(h)$ and $|B_t|=z\varphi(h)$; the latter
surjectivity follows already from the lifts just constructed.
The homomorphism $\Phi\to A\times_{U_h}B_t$ is injective, since its
kernel is $T\cap T'$, and both sides have this same finite order.
It is an isomorphism.  In the kernel over $U_h$, $t$ centralizes
$T'^c$, and their intersection is trivial; this proves the asserted
direct product $\mu_h\times Z$.

A section of the fibre product projects to a section of each factor.
Conversely a pair of sections gives a section of the fibre product.
It remains to compute the $B_t$ factor.  The following coordinate
computations take place in the rational cocharacter space modulo the
coweight lattice $P^\vee$, which is the torus of the adjoint group.

In $A_{m-1}$, the eigenvalues of a representative of $t$ are the $m$th
roots of unity, each once, with an irrelevant common scalar.  Thus
$B_t$ is the affine permutation group of $\mathbb Z/m\mathbb Z$;
the multiplicative permutations split it.
In $B_r$, $P^\vee=\mathbb Z^r$ and
$\rho^\vee=(r,r-1,\ldots,1)$, with $h=2r$.
Multiplication by $a\in U_{2r}$ permutes the pairs $\{j,-j\}$ for
$1\le j<r$; fix the coordinate of residue $r$ with positive sign.
This defines signed permutations multiplicatively and gives a section.
The remaining sign on residue $r$ is the kernel of order two.

In $C_r$, $P^\vee=\mathbb Z^r\cup(\mathbb Z+\tfrac12)^r$ and
$2\rho^\vee=(2r-1,2r-3,\ldots,1)$.  A signed permutation is in
$B_t$ precisely when it multiplies the odd residues modulo $4r$ by
a unit $u$, where its image in $U_{2r}$ is $u\bmod2r$.  Indeed the
only extra equality allowed modulo $hP^\vee$ is simultaneous addition
of $2r$ to the doubled coordinates, which on odd residues is
multiplication by $1+2r$.  Every unit modulo $4r$ permutes the pairs
of odd residues, and the residue $1$ proves faithfulness.  Hence
\[
 B_t\simeq U_{4r},\qquad B_t\to U_h\simeq U_{2r}
 \quad\text{is reduction modulo }2r.
\]
Write $2r=2^s q$ with $q$ odd.  The Chinese remainder isomorphisms
reduce the splitting question to $U_{2^{s+1}}\to U_{2^s}$, with
identity on $U_q$.  For $s=1$ the target is trivial; for $s=2$ the
section sends $-1$ to $-1$.  For $s\ge3$, the factors generated by
$5$ have orders $2^{s-1}$ and $2^{s-2}$, respectively.  Each lift
of $5\bmod2^s$ has order $2^{s-1}$ (use
$v_2(5^{2^j}-1)=j+2$, proved by successive difference-of-squares),
so no section exists.  Thus splitting holds exactly when $4\nmid r$.
The first counterexample, $U_{16}\to U_8$, is exhaustively printed
by \texttt{small\_counterexamples.py}.

In $D_r$, put $m=r-1$.  Here
$P^\vee=\mathbb Z^r\cup(\mathbb Z+\tfrac12)^r$ and
$\rho^\vee=(m,m-1,\ldots,1,0)$, with $h=2m$.
For a section use multiplication by $a\in U_{2m}$ on the pairs
$\{j,-j\}$, $1\le j<m$, fix residue $m$ positively, and choose
the sign on residue zero to make the total sign parity even.  This
is multiplicative for the same parity-character reason used in D.002.

For $G_2$ the kernel $Z$ is trivial.  For $E_7$ it has order two;
also $-I\in W$ maps to $-1\in U_{18}$ and is central.  If $b$ lies
over the generator $5\in U_{18}$, then $b^3$ lies over $-1$, so it
is $-I$ times an element of the kernel.  Both commuting factors have
order at most two; hence $b^6=1$, and $b$ gives a section.
The other exceptional types already have a nonsplit $A$ factor by
D.002.  This proves the full table.
\end{proof}

\begin{saresult}{D.004}{Proved under stated hypotheses}{The geometric count of Kostant rigidifications}
Let $G/k$ be as above, and let $\operatorname{RigKos}(k)$ be the set of
quadruples $(T,T',B,B')$ with $(T,T')$ in apposition and Borels
$B\supset T$, $B'\supset T'$.  The conjugation action of $G(k)$ is
free, and the number of its orbits is
\[
 \frac{|W|^2}{h z\varphi(h)}.
\]
If in addition a generator of the cyclic subgroup $T\cap N_G(T')$
is marked, the number is $|W|^2/(hz)$.
\end{saresult}
\begin{proof}
All pairs in apposition form one $G$-orbit by the definition and the
cyclic-element theorem, with stabilizer $\Phi$.  Fix such a pair.
The finite sets $\mathcal Q(T)=\{B:B\supset T\}$ and
$\mathcal Q(T')$ are simply transitive Weyl sets, each of cardinality
$|W|$.  The stabilizer in $\Phi$ of a point of their product lies in
$T\cap T'=1$, so the action is free.  Orbits of quadruples are exactly
orbits of this finite free action.  D.003 gives $|\Phi|=hz\varphi(h)$,
proving the count.  Marking a generator replaces $\Phi$ by the kernel
of its exponent map, of order $hz$.
All maps used here are the conjugation maps of algebraic homogeneous
spaces and their finite Weyl sets.  Since these finite groups are
reduced in characteristic zero, the corresponding quotients are
finite etale schemes, and the action on the quadruple scheme is a
$G$-torsor over that finite quotient.  For any characteristic-zero
base $S$ and a form of this group, etale descent of the split pinned
construction gives the same assertion and the same constant geometric
rank.  This last argument asserts a finite etale \emph{form}, not
that the finite etale scheme is split over $S$.
For reference the exceptional ranks are
\[
\begin{array}{c|r}
G_2&12\\ F_4&27\,648\\ E_6&18\,662\,400\\
E_7&39\,016\,857\,600\\ E_8&2\,022\,633\,897\,984\,000.
\end{array}
\]
For the classical families substitute
$|W(A_r)|=(r+1)!$, $|W(B_r)|=|W(C_r)|=2^r r!$,
and $|W(D_r)|=2^{r-1}r!$ in the displayed formula.
The script \texttt{code/rigkos\_counts.py} records exact
integer arithmetic, including the correction to the
earlier study's $E_8$ value, which omitted one factor
$|W|$.
\end{proof}

\subsection{Outer signs and the meaning of a canonical line}
\begin{saresult}{D.005}{Proved}{The highest-root and cyclic-quotient characters}
For a pinned simple group in characteristic zero the character of its
diagram automorphism group on $\mathfrak g_\theta$ equals its character
on the Kostant line $\Delta$.  This character is nontrivial exactly in
type $A_{2m}$, where diagram reversal acts by $-1$.  It is trivial in
$A_{2m-1}$, in $D_r$ (including the whole triality group for $D_4$),
and in $E_6$; all remaining simple types have trivial diagram
automorphism group.
There is, however, no isomorphism
$\mathfrak g_\theta\simeq\Delta$ natural in just $(G,T,B)$.
\end{saresult}
\begin{proof}
In $\mathfrak{sl}_{r+1}$ a pinned reversal is
$X\mapsto-JX^{\mathsf t}J^{-1}$ with an alternating anti-diagonal
matrix $J$.  It sends $E_{1,r+1}$ to
$(-1)^{r-1}E_{1,r+1}$, giving the stated parity rule for every $r$.
For $D_r$, realize the diagram involution by an orthogonal reflection
on the last coordinate.  It exchanges the two end simple roots, fixes
the other pinned simple root vectors, and fixes the highest root
$e_1+e_2$ and its vector.  Its highest-root sign is therefore $+1$.
For $D_4$ the three transpositions are conjugate in the diagram group
$S_3$, so every transposition has sign $+1$ and the whole character is
trivial.

For $E_6$ the exact root calculation is attached as
\texttt{pinned\_signs.py}.  To make its content transparent, in a
simply laced root lattice define a bimultiplicative sign $\epsilon$
by $\epsilon(\alpha_i,\alpha_i)=-1$,
$\epsilon(\alpha_i,\alpha_j)=-1$ when $i>j$ and the vertices are
adjacent, and $+1$ otherwise.  In the root-lattice construction of
the Chevalley algebra, $[e_\alpha,e_\beta]=
\epsilon(\alpha,\beta)e_{\alpha+\beta}$ when the sum is a root.
The pinned lift of a diagram permutation $\sigma$ has signs
$\eta(\alpha_i)=1$, determined recursively by
\[
\eta(\beta+\alpha_i)=\eta(\beta)
 \epsilon(\sigma\beta,\sigma\alpha_i)\epsilon(\beta,\alpha_i).
\]
The program checks independence of \emph{every} available simple-root
decomposition and checks all 1,440 nonzero root-sum brackets of $E_6$.
For the Bourbaki permutation $(1\ 6)(3\ 5)$ it obtains
$\theta=(1,2,2,3,2,1)$ and $\eta(\theta)=+1$.
The same program prints independent finite checks of the matrix formulas
in several classical ranks and of both generators of triality; the
uniform classical arguments above, not those examples, prove the
infinite families.

Here is a canonical description that compares the two characters
without conflating their underlying constructions.  Define the line
\[
 Q=\mathfrak g_{-\theta}\otimes
       \bigotimes_{i=1}^r\mathfrak g_{\alpha_i}^{\otimes m_i}.
\]
Its $T$-character is zero.  The map from tuples of nonzero root
components to $Q^\times$ sends $(x_0,\ldots,x_r)$ to
$x_0\otimes\bigotimes_i x_i^{\otimes m_i}$.  It is a $T$-invariant
surjection, and its fibres are precisely $T$-orbits: the simple root
characters, a basis for the adjoint character lattice, uniquely match
the last $r$ components, and equality of tensors then matches the
first.  Therefore the geometric cyclic quotient is canonically
$Q^\times$.  Kostant's cyclic quotient identifies it with
$\Delta^\times$; extending across zero gives an isomorphism of lines
$Q\simeq\Delta$.  A pinning trivializes the simple-root factors,
so the diagram character is that on $\mathfrak g_{-\theta}$.
The invariant nondegenerate pairing between opposite root lines makes
this character the inverse of that on $\mathfrak g_\theta$; the
characters just calculated take only the values $\pm1$, hence they
are equal.

Finally, any isomorphism natural in $(G,T,B)$ must commute with its
automorphisms given by conjugation by elements of $T$.  They act on
$\mathfrak g_\theta$ by the nontrivial root character $\theta$, and
act trivially on the adjoint-quotient line $\Delta$.  A nonzero linear
map cannot intertwine these two actions.  Thus the unqualified
canonical isomorphism asked for at 123/034 is false.  Equality of
outer characters is a strictly weaker assertion, and the tensor
line $Q$ states precisely the corrected functorial object.
\end{proof}

\subsection{Deleting one cyclic coefficient}
\begin{saresult}{D.006}{Proved}{The single-vanishing criterion in every simple type}
Let $0\ne x_j\in\mathfrak g_{\alpha_j}$ for $j\ne i$, and put
$x^{(i)}=\sum_{j\ne i}x_j$.  Then $x^{(i)}$ is nilpotent.  It is
regular nilpotent if and only if $m_i=1$.  With Bourbaki numbering the
complete list of permitted missing nodes is
\[
\begin{array}{c|l}
A_r&0,1,\ldots,r\\
B_r&0,1\\ C_r&0,r\\ D_r&0,1,r-1,r\\
E_6&0,1,6\\ E_7&0,7\\ E_8,F_4,G_2&0.
\end{array}
\]
\end{saresult}
\begin{proof}
Any $r$ of the affine simple roots $\alpha_0,\ldots,\alpha_r$ are
linearly independent: their only dependence is
$\sum_{j=0}^r m_j\alpha_j=0$, all coefficients nonzero.  There is
therefore a rational cocharacter taking value one on all the retained
roots; multiplying it by an integer gives an integral cocharacter
$\lambda:\mathbb G_m\to T$ with positive values on them.  Thus
$\lim_{s\to0}\operatorname{Ad}(\lambda(s))x^{(i)}=0$ and
$x^{(i)}$ is nilpotent.

If $m_i=1$, replacing $\alpha_i$ by $-\theta$ gives a simple system
of the full root system.  Here is an affine-arrangement verification.
The vertex opposite the $i$th wall of the fundamental affine alcove is
$v=\omega_i^\vee/m_i=\omega_i^\vee\in P^\vee$.
Translation by $-v$ preserves every affine root hyperplane and moves
this vertex to zero.  In a sufficiently small neighbourhood of zero,
the only affine root hyperplanes present are the ordinary root
hyperplanes through zero.  The local cone of the translated alcove
is consequently one full Weyl chamber; its inward wall normals are
exactly the retained roots.  For $i=0$ use $v=0$.
The sum of nonzero vectors for
a full simple system is regular nilpotent, and a torus conjugation
makes all its coefficients those of the usual principal element.

If $m_i>1$, take $s\in T$ with $\alpha_i(s)$ a primitive $m_i$th
root of unity and all other simple-root characters equal to one.
This exists uniquely because $G$ is adjoint.  It is nonidentity and
semisimple, and $(-\theta)(s)=1$, so it centralizes $x^{(i)}$.
The centralizer of a regular nilpotent element in an adjoint group
has no nonidentity semisimple element: such an element must preserve
the unique Borel containing the regular nilpotent; conjugate its
semisimple part into the torus of that Borel, where the nonzero simple
root coefficients force every simple-root character to be one.
This contradicts regularity.  The list is obtained by reading the
coefficients of the highest roots; the exceptional coefficients are
also regenerated from the Cartan matrices in the attached code.
\end{proof}

\subsection{Dynkin characteristics and parabolics}
For a nilpotent element in characteristic zero choose an
$\mathfrak{sl}_2$-triple $(e,H,f)$ with $H$ dominant.  Its weighted
Dynkin diagram is $d_i=\alpha_i(H)\in\{0,1,2\}$, and its associated
parabolic has Lie algebra $\bigoplus_{j\ge0}\mathfrak g(j)$, where
$\mathfrak g(j)=\{X:[H,X]=jX\}$.  Thus its standard parabolic class
is specified by $\{i:d_i=0\}$.  This defines the map of sets from
nilpotent adjoint orbits to conjugacy classes of parabolic subgroups.
The weighted diagram determines the orbit; for the conventions and
the exact triple criterion see \cite[\S\S2--3, Lemma 2]{D-deGraaf}.

\begin{saresult}{D.007}{Proved}{Exactly when the Dynkin-to-parabolic map is injective}
For a simple complex Lie algebra, including the zero orbit, the map
just defined is injective exactly in types $A_1$ and $B_2=C_2$.
In particular $A_2$, not $A_3$, already supplies a counterexample.
\end{saresult}
\begin{proof}
For $A_1$ the diagrams are $(0)$ and $(2)$.
For $C_2$ the four symplectic partitions of four are
$(1^4),(2,1,1),(2,2),(4)$; their diagrams are respectively
$(0,0),(1,0),(0,2),(2,2)$, with four different zero sets.

For every simply laced irreducible root system of rank at least two,
the highest coroot $\theta^\vee$ is dominant and its labels are zero
or one.  There are roots $\beta,\gamma$ with
$\theta=\beta+\gamma$ and $\langle\beta,\gamma^\vee\rangle=-1$:
choose a simple root $\alpha$ with
$\langle\theta,\alpha^\vee\rangle=1$ and put
$\beta=\alpha$, $\gamma=\theta-\alpha$.
The root $\mathfrak{sl}_2$ for $\theta$ has characteristic
$\theta^\vee$.  The principal triple of the $A_2$ subsystem generated
by $\beta,\gamma$ has characteristic $2(\beta^\vee+\gamma^\vee)
=2\theta^\vee$.  These are distinct dominant diagrams with the same
zero set, proving noninjectivity for $A_r$ ($r\ge2$), $D_r$, and
$E_6,E_7,E_8$.

The same construction works with long roots in $B_r$ ($r\ge3$),
$F_4$, and $G_2$: their long roots contain an $A_2$ subsystem with
highest root $\theta$, and $\theta^\vee$ has only labels zero and
one.  In $B_r$ one can take
$\beta=e_1-e_3$, $\gamma=e_2+e_3$; the long-root systems of $F_4$
and $G_2$ are $D_4$ and $A_2$, respectively.

In $C_r$ for $r\ge3$, the symplectic partitions
$(3,3,1^{2r-6})$ and $(2,2,1^{2r-4})$ are both allowed, since every
odd part occurs with even multiplicity.  Their dominant
characteristics in the standard coordinates are
$(2,2,0,\ldots,0)$ and $(1,1,0,\ldots,0)$.
Their simple-root labels are respectively two and one at vertex two
and zero elsewhere, so the zero sets agree.  This completes all
simple types.  The $A_2$ case explicitly has diagrams $(1,1)$ and
$(2,2)$.
\end{proof}

\begin{saresult}{D.008}{Proved}{All nilpotent classes whose Dynkin parabolic is a Borel}
In $A_r$, use partitions of $r+1$.  The classes in question are exactly
the one-part partition and the two-part partitions whose two parts
have opposite parity.  In every other simple type the only such
class is the regular nilpotent class.
\end{saresult}
\begin{proof}
For $A_r$, the weights of $H$ contributed by a Jordan block of size
$a$ in the standard representation are
$a-1,a-3,\ldots,-a+1$.  The associated parabolic is a Borel exactly
when all weights of the standard representation are distinct.
Two blocks of the same parity have an overlapping weight, whereas
two blocks of opposite parity have disjoint weight sets.  Three
blocks necessarily repeat a parity.  This proves exactly the stated
partition criterion.

For $B_r$ the coordinates of a regular Cartan element must have
distinct nonzero absolute values.  An even Jordan block has even
multiplicity in an orthogonal partition, which would repeat a
nonzero weight; two odd blocks would give at least two zero weights
and hence a zero coordinate beyond the one compulsory zero weight
in dimension $2r+1$.  Thus the only possibility is the one odd block
$(2r+1)$, the regular orbit.  For $C_r$, an odd block has even
multiplicity and supplies repeated zero weights, forbidden by
regularity.  Two even blocks overlap in the weights $1,-1$.
Consequently only $(2r)$ is possible.  For $D_r$, distinct nonzero
absolute values are required, but one zero coordinate is allowed.
Even blocks again repeat nonzero weights.  There must be precisely
two odd blocks; two of size at least three repeat $2,-2$, so one
must be of size one.  The partition is $(2r-1,1)$, again regular.

For the exceptional groups we give an exhaustive rational certificate
with a criterion special to this question.  A Borel diagram has
$d\in\{1,2\}^r$, so $\mathfrak g(0)=\mathfrak t$.  Let
$R_2(d)=\{\beta>0:\sum_i\beta_i d_i=2\}$, and let $H_d$ be the
unique Cartan vector with $\alpha_i(H_d)=d_i$.  Then
\[
 d\text{ is a nilpotent characteristic}
 \quad\Longleftrightarrow\quad
 H_d\in\operatorname{span}_{\mathbb Q}
          \{\beta^\vee:\beta\in R_2(d)\}.                 \tag{D.8.1}
\]
Indeed, for different $\beta,\gamma\in R_2(d)$ the difference has
degree zero and cannot be a root, since every positive root has
strictly positive degree.  Hence all cross brackets
$[e_\beta,f_\gamma]$ vanish.  An existing triple forces its Cartan
bracket $H_d$ into the displayed span.  Conversely write
$H_d=\sum b_\beta\beta^\vee$, normalize
$[e_\beta,f_\beta]=\beta^\vee$, and set
$e=\sum e_\beta$, $f=\sum b_\beta f_\beta$.
The cross-bracket observation gives $[e,f]=H_d$, while their degrees
give $[H_d,e]=2e$ and $[H_d,f]=-2f$.  This is the required triple.

\texttt{borel\_dynkin.py} generates every root from each exceptional
Cartan matrix, tests (D.8.1) for \emph{all} $2^r$ positive diagrams,
and records either rational coefficients $b_\beta$ or a rational
linear functional annihilating every $\beta^\vee$ but not $H_d$.
The complete outputs are
\[
\begin{array}{c|r|l}
 G_2&4&(2,2)\\ F_4&16&(2,2,2,2)\\
 E_6&64&(2,2,2,2,2,2)\\
 E_7&128&(2,2,2,2,2,2,2)\\
 E_8&256&(2,2,2,2,2,2,2,2).
\end{array}
\]
The middle column is the number of tested diagrams and the last
column is the complete accepted list.  This proves the exceptional
assertion without assuming that a random unsuccessful search for
an $\mathfrak{sl}_2$-triple certifies nonexistence.
\end{proof}

\subsection{A precise relative theorem in characteristic zero}
\begin{saresult}{D.009}{Proved under stated hypotheses}{Relative simultaneous resolution}
Let $S$ be a scheme over $\mathbb Q$.  In the Lie-algebra version let
$G/S$ be semisimple, with any isogeny type; in the group version let
$G/S$ be reductive with simply connected derived subgroup.  Let $D/S$
be its flag scheme and $T/S$ its torus determined by the based root
datum, and let $W/S$ be the corresponding finite etale Weyl group.
Put $(X,I')=(\mathfrak g,\mathfrak t)$ in the Lie-algebra version and
$(X,I')=(G,T)$ in the group version, and set $I=I'/W$.
There is a canonical adjoint quotient $f:X\to I$ and a diagram
\[
 \widetilde X\xrightarrow{p}X'=X\times_I I'
   \xrightarrow{q}X,
 \qquad \widetilde f:\widetilde X\to I',
\]
where $\widetilde X$ is the universal Borel Lie algebra or the universal
Borel group over $D$, respectively.  The maps $p$ and $q$ here have
the indicated scheme-theoretic targets; $q$ is finite and $p$ is
projective.  The schemes $I$ and $I'$ are smooth over $S$, $f$ is flat
with geometrically normal fibres, and $\widetilde f$ is smooth.
The regular locus $X^{\rm reg}$ is the smooth locus of $f$.
The map $p$ is an isomorphism over $X'^{\rm reg}$, and no larger
open subset has this property.  Fibrewise over $S$, the complements
of the regular loci in $\widetilde X$ and $X'$ have codimension at
least two and pure codimension three, respectively.  If $S$ is normal
and locally noetherian, then $X'$ is normal.
\end{saresult}
\begin{proof}
We specify both the field input and why it transfers to the asserted
base; normality of the generic fibre alone would not suffice.
Over an algebraically closed characteristic-zero field the adjoint
quotient, its normal fibres, its regular smooth locus, and its
simultaneous resolution are the Kostant--Steinberg theorems.  For the
complex group formulation the simultaneous-resolution diagram and
regularity assertions are explicitly stated by Brieskorn
\cite[pp.279--280]{D-Brieskorn}; the unipotent restriction and its
separability hypotheses are stated by Slodowy
\cite[\S4.1, pp.43--44]{D-Slodowy}.  The codimension assertions are
the regular/subregular part of this field theorem.  These field
assertions are used as known results, not as newly proved statements.

First take a split semisimple simply connected group over
$\mathbb Q$.  Its group quotient is affine space, with coordinates
the fundamental characters; its Lie quotient is affine space with
the fundamental homogeneous invariant polynomials.  The restriction
maps identify these rings with the Weyl invariants on $T$ or
$\mathfrak t$.  Over $\mathbb Q$ the invariant projector is a
direct-summand projector, so these invariant-ring constructions
commute with every extension of the coefficient algebra, including
nonflat extensions.  Thus the entire diagram over any
$\mathbb Q$-scheme is the base change of the diagram over
$\mathbb Q$.  The complex field assertions descend faithfully flatly
to $\mathbb Q$ and then ascend by base change: flatness, smoothness,
projectivity and geometrically normal fibres have exactly these
descent and base-change properties.  Fibrewise codimension is checked
after algebraic closure of each residue field, where the field theorem
applies.  Central isogenies do not change semisimple Lie algebras in
characteristic zero, giving the Lie version for every isogeny type.

For a split reductive group with simply connected derived subgroup,
the central isogeny
$G_{\rm der}\times\operatorname{Rad}(G)\to G$ has a finite kernel $K$ whose
projection to the central torus is injective.  On the quotient
$\mathbb A^r\times\operatorname{Rad}(G)$ its action is therefore free.  Both
the group and quotient maps are $K$-torsors, and their square is
cartesian: a lift of the group element and a lift of its quotient
value differ by exactly one element of $K$.
Consequently the map on quotients is finite etale, and all the above
properties descend from the product of the semisimple diagram with
the central torus.  This proves the split reductive version.

For a form over $S$, choose an etale cover on which the group is
split with fixed root datum.  The incidence scheme and the quotient
are equivariant under all automorphisms of the pinned split model;
inner automorphisms act trivially on the quotient.  The local
diagrams hence have canonical descent data and descend to the stated
diagram over $S$.  All the asserted properties are etale local.
In particular this reduction is through actual comparison
isomorphisms, not an appeal to a generic fibre of an integral model.

One can also see smoothness of $\widetilde f$ directly: locally the
map to $D\times_S I'$ is a torsor for a smooth unipotent group,
followed by the smooth projection $D\times_S I'\to I'$.
The incidence embedding in $D\times_S X$ is closed, which proves
projectivity of $p$.  Over the regular locus it is proper with finite
fibres, hence finite.  On every geometric fibre it is an isomorphism
by the field theorem.  Both its source and target are flat over $S$;
the finite algebra map is therefore an isomorphism by Nakayama's
lemma applied first to its cokernel and then, using flatness, to its
kernel.  Outside this locus the geometric fibre has positive
dimension, so an isomorphism is impossible.  Finally a flat morphism
with geometrically normal fibres over a normal base has normal total
space, which applies to $X'\to I'$.
\end{proof}

\begin{saresult}{D.010}{Proved under stated hypotheses}{Picard sequences, including the character term}
Use the group version of D.009 and assume that $S$ is connected,
normal, and locally noetherian.  Write $\xi$ for its generic point,
$M=X^*(T)$ for the etale character-lattice sheaf, and
$\pi:\mathcal B=\widetilde X\to D$ for the universal Borel group.
Then pullback and the unit section give canonical decompositions
\[
 \operatorname{Pic}(\mathcal B)
   \simeq\operatorname{Pic}(D)\oplus H^1_{\rm et}(S,M),
 \qquad
 \operatorname{Pic}(T)
   \simeq\operatorname{Pic}(S)\oplus H^1_{\rm et}(S,M),
                                                        \tag{D.10.1}
\]
and restriction and $p$ induce isomorphisms
\[
 \operatorname{Pic}(\mathcal B)
 \simeq\operatorname{Pic}(\mathcal B^{\rm reg})
 \simeq\operatorname{Pic}(G'^{\rm reg}).                   \tag{D.10.2}
\]
Pullback along $G'\to T$ is an isomorphism
$\operatorname{Pic}(T)\simeq\operatorname{Pic}(G')$.
In particular the manuscript's sequences hold, and in this scope can
be completed by zeros on the right:
\[
\begin{split}
0&\to\operatorname{Pic}(D)\to\operatorname{Pic}(\mathcal B)
       \to H^1_{\rm et}(S,M)\to0,\\
0&\to\operatorname{Pic}(S)\to\operatorname{Pic}(G')
       \to H^1_{\rm et}(S,M)\to0.
\end{split}                                             \tag{D.10.3}
\]
If $\Gamma$ is the finite monodromy group of $M$, then
$H^1_{\rm et}(S,M)=H^1(\Gamma,M_{\bar\xi})
=H^1(\xi,M)$, a finite group.  After etale sheafification this term
vanishes, giving the relative Picard sheaf isomorphisms of 123/011.
Finally
\[
0\to\operatorname{Pic}(S)\to\operatorname{Pic}(G)
       \to\operatorname{Pic}(G_\xi)
\]
is exact, and restriction
$\operatorname{Pic}(G)\to\operatorname{Pic}(G^{\rm reg})$
is an isomorphism.  The group $\operatorname{Pic}(G_\xi)$ is finite;
it is zero when $G$ is semisimple simply connected, even for a
nonsplit form.
\end{saresult}
\begin{proof}
Here is a divisor argument that handles the normal, not necessarily
regular, base.  If $A$ is a normal noetherian domain and
$R=A[x_1,\ldots,x_a,t_1^{\pm1},\ldots,t_b^{\pm1}]$, then
\[
 \operatorname{Cl}(A)\xrightarrow{\sim}\operatorname{Cl}(R).
                                                        \tag{D.10.4}
\]
To prove surjectivity, restrict a Weil divisor to
$\operatorname{Frac}(A)[x,t^{\pm1}]$, a unique factorization domain,
and subtract a principal divisor eliminating all horizontal
components.  Every remaining height-one prime is the extension of
a height-one prime of $A$, with multiplicity one.  For injectivity,
a rational function with no horizontal zero or pole is a constant
from $\operatorname{Frac}(A)^\times$ times a Laurent monomial; its
vertical divisor is principal already on $A$.
If $U\subset\operatorname{Spec}R$ has complement of codimension at
least two in every geometric fibre over $A$, it contains all
height-one points of $\operatorname{Spec}R$, and
$\operatorname{Cl}(U)=\operatorname{Cl}(R)$.  Moreover $U\to
\operatorname{Spec}A$ is faithfully flat.  The pullback of a
rank-one reflexive module $N\subset K=\operatorname{Frac}(A)$
representing a class of $A$ is the module representing the
pulled-back divisor.  Here is the required reflexivity check,
which must not be inferred merely from agreement at height one.
Write $N=\bigcap_{\operatorname{ht}\mathfrak q=1}N_{\mathfrak q}
\subset K$.  The horizontal height-one conditions in $R$ force a
rational element of the pulled-back divisorial module to lie in
$K[x,t^{\pm1}]$.  The vertical conditions say that each of its
finitely many coefficients lies in every $N_{\mathfrak q}$.
Their intersection is consequently
$N[x,t^{\pm1}]=N\otimes_A R$, proving that the tensor pullback
itself is the divisorial, hence reflexive, module.  If its
restriction to $U$ is invertible,
faithfully flat descent of finite locally free modules shows that
the original module is invertible.  Therefore
\[
 \operatorname{Pic}(A)\simeq\operatorname{Pic}(R)
                 \simeq\operatorname{Pic}(U).            \tag{D.10.5}
\]
Regular functions and units likewise extend across the omitted
codimension-two set by normality.

Etale locally on $D$, choose a maximal torus of its Borel group,
split it, and order the positive root groups.  The underlying variety
is then an affine space times a split torus.  Apply (D.10.5) over
each normal affine chart of $D$; the same applies to its regular
open subset by D.009.  Thus the sheaf of relative units and the
first higher direct image on the small etale site of $D$ are
\[
 \pi_*\mathbb G_m=\mathbb G_m\oplus M_D,
 \qquad R^1\pi_*\mathbb G_m=0.                           \tag{D.10.6}
\]
The splitting of units is canonical: divide a unit by its value on
the identity section; what remains is the uniquely determined
character.  The same formulas hold after removing the nonregular
locus.  The degree-one Leray sequence consequently gives
$\operatorname{Pic}(\mathcal B)=\operatorname{Pic}(D)
\oplus H^1(D,M_D)$, and proves (D.10.2).

For completeness, the identification of the last cohomology group
does not require an unmentioned rational point on the flag scheme.
Take the finite connected etale cover $S_1\to S$ with group
$\Gamma$ trivializing $M$.  Such a cover exists also when $S$
is only locally noetherian: the finite splitting extension of the
generic lattice gives its normalization, and etale local
triviality of the lattice identifies this normalization locally
with a finite etale cover.  Its pullback $D_1$ is connected and
normal because flag fibres are geometrically integral.  For any
connected normal locally noetherian scheme $Y$, one has
$H^1_{\rm et}(Y,\mathbb Z)=0$: a torsor restricts at the generic
point to a continuous homomorphism from a profinite Galois group to
$\mathbb Z$, which is zero, and a generic section extends uniquely
on each connected normal etale trivializing chart, hence globally.
The same argument works for $\mathbb Z^r$.  The covering spectral
sequence for $S_1/S$ and $D_1/D$ now gives, respectively,
$H^1(S,M)=H^1(\Gamma,M_{\bar\xi})=H^1(D,M_D)$.
Applying the identical argument at $\xi$ gives $H^1(\xi,M)$.
For a finite group acting on a finitely generated lattice, $H^1$
is finitely generated and annihilated by $|\Gamma|$ (average a
cocycle over the group), so it is finite.  It vanishes on the
trivializing etale cover, explaining precisely the relative-sheaf
assertion.  Applying (D.10.6) to $T\to S$ proves the second
decomposition in (D.10.1).

Since $p$ is proper birational and $G'$ is normal, $p_*\mathcal
O_{\mathcal B}=\mathcal O_{G'}$; projection formula makes
$p^*:\operatorname{Pic}(G')\to\operatorname{Pic}(\mathcal B)$
injective.  The identity section $D\to\mathcal B$ maps in $G'$ to
the identity section $S\to G'$.  Hence the $\operatorname{Pic}(D)$
component of any class in the image is pulled back from
$\operatorname{Pic}(S)$.  Conversely every such component and
every $H^1(S,M)$ component is the pullback of a line bundle on $T$
by (D.10.1).  Injectivity of $p^*$ then proves
$\operatorname{Pic}(G')=\operatorname{Pic}(T)$, as claimed.

For the sequence concerning $G$, a line bundle trivial on $G_\xi$
is represented, after subtracting a rational principal divisor, by
a vertical divisor on $G$.  Smoothness and geometrically integral
fibres show that every vertical prime divisor is the pullback of
a prime divisor of $S$ with multiplicity one.  Faithfully flat
descent, as in (D.10.5), shows that this divisor on $S$ is Cartier.
The unit section proves injectivity from $\operatorname{Pic}(S)$.
This is also precisely the application of EGA IV, Corollary
21.4.13 in Errata et Addenda, liste 3, p.361: its hypotheses,
specified in 21.4.9 and 21.4.11, are faithful flatness, openness
or quasi-compactness, a normal locally noetherian base, and
integral fibres over its height-one points.  Smoothness and
geometrically integral group fibres verify each hypothesis
\cite[pp.358--361]{D-EGAIV4}.
For restriction to $G^{\rm reg}$, cover the split group etale locally
by translates of the big cell, each an affine space times a torus;
translations through geometric points lift after an etale change
of the base by smoothness.  Apply (D.10.5) on these charts and glue
their unique extensions, using normality.  This gives the asserted
isomorphism.

Over an algebraic closure the Picard group of a reductive group is
finite: the complement of its big cell has finitely many prime
divisors, and a positive multiple of each is the divisor of a
fundamental highest-weight matrix coefficient.  If the derived
subgroup is simply connected the fundamental weights themselves
extend after a character of the central torus is chosen, so each
such divisor is principal and the geometric Picard group is zero.
Galois descent then injects the kernel of
$\operatorname{Pic}(G_\xi)\to\operatorname{Pic}(G_{\bar\xi})$
into $H^1(\xi,X^*(G_{\bar\xi}))$; the scalar-unit contribution is
zero by Hilbert 90.  This lattice cohomology is finite by the
argument just given.  For semisimple simply connected $G$, both
the geometric Picard group and the character lattice vanish,
proving the final assertion.

The ampleness assertion on 123/011--013 needs no normality:
the morphism $\pi:\mathcal B\to D$ is affine and has a closed
unit section, so for a line bundle $L$ on $D$, $L$ is ample
relative to $S$ if and only if $\pi^*L$ is.  One direction is
pullback under an affine morphism, and the converse is
restriction to the closed unit section.  On the flag variety,
with the sign convention giving dominant weights nef line
bundles, the ample weights are the \emph{strictly} dominant
ones.  Thus the manuscript's notation $P^+$ denotes the ample
cone only if it means the open dominant chamber, not its usual
closed dominant cone.
\end{proof}

\begin{saresult}{D.011}{Proved under stated hypotheses}{The projectivity obstruction on 123/014}
Let $G$ be a connected reductive group with nontrivial semisimple
part over an algebraically closed characteristic-zero field.
There is no projective birational morphism $v:Y\to G$ such that
$Y$ is smooth, $v$ is an isomorphism over $G^{\rm reg}$,
$Y\setminus v^{-1}(G^{\rm reg})$ has codimension at least two,
and the composite $Y\to G\to G/\!/G$ is smooth.
If the derived subgroup of $G$ is simply connected, the
codimension-at-least-two condition follows from the other
conditions and need not be imposed separately.
\end{saresult}
\begin{proof}
The smooth group $G$ has finite Picard group, by the last argument
of D.010.  Removing a codimension-at-least-two subset from a regular
scheme does not change its Picard group (Cartier and Weil divisors
coincide), so the assumptions imply
$\operatorname{Pic}(Y)\simeq\operatorname{Pic}(G^{\rm reg})
\simeq\operatorname{Pic}(G)$, a finite group.
If every fibre of $v$ were finite, properness would make $v$ finite,
and normality of $G$ would make the birational map an isomorphism.
But the adjoint quotient is not smooth at a nonregular unipotent
element, so the last assumption rules this out.  There is therefore
a positive-dimensional proper fibre $F$.  A relatively ample line
bundle for the projective morphism $v$ restricts to an ample bundle
on $F$, while its class is torsion because $\operatorname{Pic}(Y)$
is finite.  An ample torsion line bundle would make
$\mathcal O_F$ ample, forcing the proper scheme $F$ to be affine
and zero-dimensional.  This is a contradiction.  The same argument
on a geometric fibre gives the relative obstruction whenever all
four assumptions hold fibrewise.

For the additional assertion, write $I=G/\!/G$, $d=\dim G-r$,
and suppose that $G_{\rm der}$ is simply connected.
The quotient fibres $G_a$, $a\in I$, are integral, and
$G_a^{\rm reg}$ is dense and nonempty, by D.017.  Proper
birationality over the normal group gives connected fibres
for $v$ by Stein factorization.  Each fibre $Y_a$ of the
smooth composite is smooth of pure dimension $d$.  Its
component containing the irreducible open
$v^{-1}(G_a^{\rm reg})$ maps properly and surjectively to
$G_a$.  Any other component would be disjoint from this
one and would give a disconnected fibre of $v$ over each
point in its nonempty image.  Consequently $Y_a$ is
irreducible.  This argument works over an algebraic
closure of every residue field of $I$.

If an exceptional prime divisor $E$ existed, its image
in $I$ would lie in the discriminant, since every element
over the regular semisimple open set is regular.  Let
$Z$ be the closure of this image.  Since
$\dim E=\dim G-1$ and $\dim Z\le r-1$, the generic
fibre of $E\to Z$ has dimension at least $d$.  The
smoothness of $Y\to I$ bounds it above by $d$, so it
is an entire component of the geometrically irreducible
fibre $Y_\eta$.  It therefore meets
$v^{-1}(G_\eta^{\rm reg})$, contrary to exceptionality.
There are no exceptional divisors.  Thus the boundary
in the statement has codimension at least two, proving
the additional assertion.
\end{proof}

\subsection{Counterexamples and other questions in the letter}
\begin{saresult}{D.012}{Refuted}{The strict differential-rank guess in nonsimply-laced type}
The letter's guess that the rank of the adjoint quotient differential
at a subregular nilpotent is strictly smaller than $r-1$ in
nonsimply-laced types is false.  It already fails in type $C_2$ over
$\mathbb C$.
\end{saresult}
\begin{proof}
In the standard symplectic realization take
$X=\left(\begin{smallmatrix}0&I_2\\0&0\end{smallmatrix}\right)
\in\mathfrak{sp}_4$.  The Lie algebra consists of matrices
$Y=\left(\begin{smallmatrix}A&B\\C&-A^{\mathsf t}\end{smallmatrix}\right)$
with $B,C$ symmetric.  Directly solving $[X,Y]=0$ gives $C=0$ and
$A=-A^{\mathsf t}$, while $B$ is arbitrary symmetric.  The
centralizer has dimension $1+3=4=r+2$, so $X$ is subregular.
The invariant algebra is generated by $\operatorname{tr}(X^2)$
and $\det X$ (equivalently by the two nonconstant characteristic
coefficients).  Their differentials at $X$ are respectively
$Y\mapsto2\operatorname{tr}C$, a nonzero functional, and zero,
since every $3\times3$ minor of $X$ vanishes.  Thus the rank is
$1=r-1$.  The complete ten-dimensional commutator and differential
calculation is in \texttt{small\_counterexamples.py}; its exact
output gives commutator rank six, centralizer dimension four, and
quotient rank one.
\end{proof}

\begin{saresult}{D.013}{Proved}{Some opposite Borels always exist; all need not be opposite}
For any two maximal tori $T,T'$ of a connected reductive group over
an algebraically closed field, and any Borel $B\supset T$, there is
a Borel $B'\supset T'$ opposite to $B$.  The stronger assertion
that every pair of Borels containing two apposed tori is opposite
is false, already for $\operatorname{PGL}_4(\mathbb C)$.
\end{saresult}
\begin{proof}
The Borels containing $T'$ form an apartment, with one chamber for
each element of the finite Weyl group.  Choose such a chamber $C$
for which the length of the relative Bruhat position $\delta(B,C)$
is maximal.  If this position $w$ is not the longest element,
there is a simple reflection $s$ with $\ell(ws)=\ell(w)+1$.
Move from $C$ to the unique $s$-adjacent chamber in that apartment.
The rank-one Bruhat multiplication rule says that its relative
position with respect to $B$ is $ws$: for adjacent chambers it
lies in $\{w,ws\}$, and when the latter has larger length that
value is forced.  This contradicts maximality.  Hence the maximal
position is the longest element, which is precisely opposition.
This argument uses the ordinary Bruhat adjacency rule, valid over
every algebraically closed field, and is independent of apposition.

For the negative assertion let $T'$ be the diagonal torus in
$\operatorname{PGL}_4$, and let $T=F T'F^{-1}$, where
$F=(i^{ab})_{0\le a,b<4}$ is the Fourier matrix.  It diagonalizes
the cyclic shift, so this is an apposed pair.  Its determinant is
$-16i$, but its minor with rows $\{0,2\}$ and columns $\{0,2\}$
is zero.  Take a coordinate flag whose two-plane is
$\langle e_1,e_3\rangle$ and a Fourier eigenflag whose two-plane
is spanned by columns zero and two of $F$.  The vanishing minor
says these two planes are not complementary, so the flags are
not opposite.  Both flags are completed by ordering their
respective eigenvectors, and their stabilizers are Borels
containing the indicated tori.  The nonzero determinant and zero
minor are checked exactly in the attached computation.
\end{proof}

\begin{saresult}{D.014}{Proved under stated hypotheses}{Rational singularities of characteristic-zero quotient fibres}
Every fibre of the adjoint quotient of a complex reductive Lie
algebra has rational singularities.  The same is true for the group
quotient when the derived group is simply connected.
\end{saresult}
\begin{proof}
The known input is Hesselink's theorem that the nilpotent variety of
a split reductive Lie algebra in characteristic zero has a rational
resolution \cite[Theorem A, p.249]{D-Hesselink}.  We explain the
passage to other fibres, which that formulation does not state.
For a fixed semisimple element $s$, let $L=C_G(s)^\circ$ and let
$\mathcal N(\mathfrak l)$ be the nilpotent cone of its Lie algebra.
The map
\[
 G\times^L\mathcal N(\mathfrak l)\longrightarrow\chi^{-1}(\chi(s)),
 \qquad[g,n]\longmapsto\operatorname{Ad}(g)(s+n)
\]
is an isomorphism.  Indeed Jordan decomposition identifies its
points uniquely modulo $L$; on this fixed fibre the semisimple part
is a polynomial in the matrix in any faithful representation, by
Hermite interpolation on the fixed eigenvalues and sufficiently
high multiplicities.  It therefore gives a regular map to the
orbit $G/L$, with fibre $\mathcal N(\mathfrak l)$, and supplies
the algebraic inverse after the smooth local trivializations of
$G\to G/L$.  Semisimple Lie centralizers are connected Levi
subgroups in characteristic zero, so no component quotient is
omitted.
The associated bundle is etale locally the product of a smooth
variety with this nilpotent cone.  Pulling back Hesselink's
resolution proves vanishing of its higher direct images locally;
flat base change and etale descent give rational singularities.
For the group version semisimple centralizers are connected under
the simply connected derived-subgroup assumption.  Replace
$s+n$ by $su$, use multiplicative Jordan decomposition, and use
the characteristic-zero polynomial exponential/logarithm
isomorphism between the unipotent variety of $L$ and
$\mathcal N(\mathfrak l)$.  The same argument applies.
This proves the rational-singularity question of the letter in
the stated characteristic-zero scope.
\end{proof}

\subsection{The integral rigidification scheme and its nonsplit form}
\begin{saresult}{D.015}{Proved under stated hypotheses}{The RigKos rank over the full base with $h$ invertible}
Let $S$ be a scheme on which the Coxeter number $h$ is invertible,
and let $G/S$ be an adjoint simple group of fixed type.  Define an
apposed pair etale locally as follows: a generator $n$ of a cyclic
group of order $h$ normalizes one maximal torus $T_0$, induces a
Coxeter element on it, and the other torus is $C_G(n)^\circ$.
This is independent of the chosen generator.  The scheme of
apposed pairs so defined agrees with the cyclic-line definition
in the notes, over the entire stated base, by D.021 below.
The scheme of
quadruples $(T,T',B,B')$ with an apposed pair and Borels containing
the respective tori is smooth over $S$, its conjugation action by
$G$ is free, and its fppf quotient is finite etale of rank
\[
 |W|^2/(h z\varphi(h)).
\]
For the split group it is a finite etale scheme over
$\mathbb Z[1/h]$; it need not be a split finite scheme.
\end{saresult}
\begin{proof}
We construct the split model over $R=\mathbb Z[1/h]$ and check every
prime that remains on this base.  Fix its split torus $T_0$ and
Chevalley representatives $n_i$ of the simple reflections, and put
$n=n_1\cdots n_r$.  In the adjoint group $n^h=1$: this is Kostant's
characteristic-zero lift theorem \cite[Theorem 8.4]{D-Kostant1959},
and the identity extends over $\mathbb Z$ because the group scheme
is separated and flat and its rational fibre is schematically dense.
Its image $c$ in $W$ has order exactly $h$, so no smaller power is
the identity in any fibre.

The centralizer $C_G(n)$ is smooth over $R$.  To see the smoothness
claim directly, infinitesimal lifts of a fixed point under the
conjugation action of the finite cyclic group can be averaged,
since its order $h$ is invertible.  Its Lie algebra has rank $r$:
on $\operatorname{Lie}T_0$ the fixed module is zero because
$\det(1-c)=z$ is invertible in $R$; every root orbit has length $h$
and contributes one fixed vector, as $n$ cyclically permutes its
root lines and the product of the coefficients is one by $n^h=1$.
There are $|R_{\rm roots}|/h=r$ such orbits.  All the primes
dividing $z$ divide $h$, as follows from
\[
(h,z)=(r+1,r+1),\ (2r,2),\ (2r,2),\ (2r-2,4),\
(12,3),(18,2),(30,1),(12,1),(6,1)
\]
in types $A,B,C,D,E_6,E_7,E_8,F_4,G_2$, respectively.
Each fibre of $n$ is semisimple, so its connected centralizer is
reductive of rank $r$ and dimension $r$, hence a torus.  Thus
$T=C_G(n)^\circ$ is a maximal torus over $R$.

We check the finite component issue rather than assuming that it
cannot change in small characteristic.  Pass to the finite etale
cyclotomic cover $R[\zeta_h]$, and set
$t=\rho^\vee(\zeta_h)\in T_0$.  In characteristic zero $n$ and $t$
are conjugate.  Their character values in every integral Weyl
module are consequently equal over this cyclotomic ring.  Weyl
characters, by their triangular highest-weight expansions, span
the invariant group ring $\mathbb Z[X^*(T_0)]^W$; finite-group
invariants separate its geometric torus orbits in every
characteristic.  Hence in every geometric fibre $n$ and $t$,
both semisimple, are conjugate.  Here one also needs a relative
conjugacy assertion, not just conjugacy on geometric fibres.
For a square-zero extension of rings, a discrepancy between two
conjugate lifts of an element of order $h$ is a $1$-cocycle of
the cyclic group on the corresponding Lie module.  Its first
cohomology is zero, by averaging with $1/h$, so the discrepancy
is removed by an infinitesimal conjugation.  Iterating this
argument proves that the transporter of $n$ to $t$ is formally
smooth.  It has finite presentation and nonempty geometric
fibres, hence is smooth and surjective, and has sections etale
locally.  This proves the required relative assertion.
No root of $t$ has value one, because its height has absolute value
less than $h$.  The root-subgroup description of its centralizer
therefore gives
\[
C_G(t)^\circ=T_0,\qquad
C_G(t)/T_0=\operatorname{Stab}_{W}(\rho^\vee/h\bmod P^\vee)
=P^\vee/Q^\vee.
\]
The last equality is the interior-alcove argument of D.003 and is
an identity of the discrete root datum, unchanged at the remaining
primes.  This component group is finite etale of rank $z$.

The fixed group $T_0^c$ is finite etale of rank $z$, since its
character group is the cokernel of $1-c$ and $z$ is invertible.
Its map to $C_G(n)/T$ is an isomorphism on the rational fibre by
D.003.  A map between finite etale group schemes is an isomorphism
on an open-and-closed subset of the connected base; consequently
it is an isomorphism everywhere.  Its kernel is $T\cap T_0$, so
this intersection is the trivial group scheme, including at every
remaining prime.

Let $\Phi=N_G(T)\cap N_G(T_0)$.  Since an element of
$T\cap N_G(T_0)$ centralizes $n$, its Weyl image centralizes $c$,
and D.002 makes that image a power of $c$.  Subtracting the
corresponding power of $n$ leaves an element of $T\cap T_0=1$.
Thus $T\cap N_G(T_0)=\langle n\rangle$.  Every element of $\Phi$
normalizes this cyclic group, and its Weyl image belongs to $A$
of D.002.  Conversely, lift any $w\in A$ with exponent $a$ to
$N_G(T_0)$.  Its conjugate of $n$ differs from $n^a$ by an element
of $T_0$.  The map $u\mapsto u c^a(u)^{-1}$ on $T_0$ is a finite
etale isogeny of degree $z$, so etale locally we can correct that
lift to conjugate $n$ exactly to $n^a$.  It then normalizes $T$.
The kernel is $T_0^c$, by the just-proved intersection formula.
Therefore $\Phi\to A_R$ is a finite etale $T_0^c$-torsor and
$|\Phi|=zh\varphi(h)$ in every geometric fibre.  This argument
does not require $\varphi(h)$ or $|W|$ to be invertible.

The Borel schemes $\mathcal Q(T)$ and $\mathcal Q(T_0)$ are finite
etale, of rank $|W|$ each.  The action of $\Phi$ on their product
is free, since a point stabilizer is $T\cap T_0=1$.  Their
quotient is therefore finite etale, with the stated rank.  The
quadruple scheme is the associated scheme
$G\times^\Phi(\mathcal Q(T)\times\mathcal Q(T_0))$; over that
finite quotient it is a $G$-torsor, hence smooth and free.
Any other Coxeter-normalized pair is etale locally conjugate to
this one: conjugate its torus to $T_0$, conjugate its Weyl Coxeter
element to $c$, and correct its lift using the same $1-c$
isogeny.  Thus this associated scheme represents exactly the
specified quadruples, not merely a selected component.
Finally every form of the split group is etale locally a split
group with the same based root datum.  The definition by torus
normalizers and neutral centralizers is invariant under all
automorphisms, so the construction and its free quotient descend
to $S$.  This proves the integral rank assertion, including
characteristics two and three whenever they do not divide $h$.
\end{proof}

\begin{saresult}{D.016}{Refuted}{Splitness of the finite quotient asked about in 123/035}
For split adjoint $A_2$ over $\mathbb Q$, the two-point geometric
quotient $\operatorname{RigKos}/G$ is
$\operatorname{Spec}\mathbb Q(\zeta_3)$.  It is not a disjoint
union of two rational points.  Thus splitness over
$\mathbb Z[1/h]$ has a negative answer.  The scan changes
``est bel et bien'' by crossing out the last three words
and inserting ``il''; it can be read as the question
``est-il splitt\'e''.  Under that reading this is an answer
to a question, not a correction of an asserted theorem;
the mathematical counterexample is the same under either reading.
\end{saresult}
\begin{proof}
Let $G=\operatorname{PGL}_3$, let $T_0$ be the diagonal torus, and
let $n$ be a three-cycle permutation matrix.  The torus
$T=C_G(n)^\circ$ and the apposed pair are defined over $\mathbb Q$.
The three eigenspaces of $n$ are defined over
$K=\mathbb Q(\zeta_3)$.  The nontrivial automorphism of $K$ fixes
the eigenline for one and exchanges the other two.  Consequently
it acts on the six Borels of $T$ by a transposition, whereas the
six Borels of $T_0$ are rational.

Here both Weyl factors are $S_3$, and D.003 identifies
\[
\Phi=\{(a,b)\in S_3\times S_3:
                    \operatorname{sgn}(a)=\operatorname{sgn}(b)\}.
\]
It has order $18$.  Its two orbits on $S_3\times S_3$ are the
two values of $\operatorname{sgn}(a)\operatorname{sgn}(b)$.
Galois acts by a transposition in exactly one factor, so exchanges
these orbits.  This is precisely the finite Galois set represented
by $\operatorname{Spec}K$.  The attached program
\texttt{rigkos\_a2.py} enumerates all 36 points, all 18 group
elements and both complete orbits, verifies the free action, and
checks Galois conjugation of the Fourier eigenbasis exactly.
Its printed Galois permutation is $[1,0]$.
The rank $36/(3\cdot3\cdot2)=2$ is correct; splitness is the
distinct assertion refuted here.
\end{proof}

\begin{saresult}{D.017}{Proved under stated hypotheses}{The group theorem over an arbitrary base, including characteristics two and three}
The group version of D.009, and consequently the normal-base
Picard assertions D.010, remain true over an arbitrary base $S$
when $G/S$ is reductive with simply connected derived subgroup.
No restriction on residual characteristic is necessary for these
group assertions.  This does not extend the adjoint Lie-algebra
assertion refuted in D.001.
\end{saresult}
\begin{proof}
We give the integral comparison, including the scheme structure
of the fibres.  In the split semisimple simply connected case the
fundamental Weyl modules over $\mathbb Z$ give regular invariant
characters $\chi_1,\ldots,\chi_r$ on $G_{\mathbb Z}$.  The
highest-weight triangular expansion proves
\[
 \mathbb Z[X^*(T)]^W=\mathbb Z[\chi_1|_T,\ldots,\chi_r|_T].
\]
The invariant torus ring commutes with every base change, because
its orbit-sum basis is a basis of a permutation module.  Thus these
characters define the universal map
$f_{\mathbb Z}:G_{\mathbb Z}\to\mathbb A^r_{\mathbb Z}=T/W$.
Restriction is also injective for invariant functions over an
arbitrary base: the conjugates of the regular semisimple torus
cover the regular semisimple locus smoothly, and that locus is
schematically dense in each smooth group fibre and in the flat
family.  Surjectivity of restriction follows from the integral
characters.  Hence this map has the intended quotient property,
not merely the intended values on field-valued points.

Over every algebraically closed field, with no restriction on its
characteristic, Steinberg proves that the set-theoretic fibres
have codimension $r$, consist of a unique dense regular class,
and have nonregular complement of codimension at least two in
the fibre \cite[Theorem 6.11, p.64]{D-Steinberg1965}.  He proves
that the differentials of the $\chi_i$ are independent exactly
on the regular locus \cite[Theorem 1.5, p.50; Theorem 8.1,
pp.70--71]{D-Steinberg1965}.  A smooth group is Cohen--Macaulay,
and the quotient is regular.  The equidimension assertion and
the local flatness criterion therefore make the quotient flat.
Each scheme-theoretic fibre is a complete intersection and is
Cohen--Macaulay.  It is smooth along its dense regular class,
hence generically reduced, hence reduced; it is regular in
codimension one and satisfies $S_2$, hence is normal.  This
supplies the scheme-theoretic strengthening needed in the
manuscript from exactly the field statements in the source.

The integral group and quotient are flat over $\mathbb Z$.
The fibrewise criterion for flatness now proves that
$f_{\mathbb Z}$ is flat; its fibres are geometrically normal by
the preceding paragraph.  The base-changed space
$G'_{\mathbb Z}=G_{\mathbb Z}\times_{\mathbb A^r}T_{\mathbb Z}$
is normal.  Its incidence morphism is projective, birational over
the regular semisimple open set, and finite over the regular
locus by Steinberg's finite-Borel criterion
\cite[Theorem 1.1, p.49]{D-Steinberg1965}.  Finite birational
morphisms to a normal scheme are isomorphisms.  Thus it is an
isomorphism over the regular locus already over $\mathbb Z$;
base change proves this over every $S$.  Outside this locus the
Borel fibre has positive dimension by the same criterion, which
also proves maximality.

Steinberg's Theorem 1.3(b), p.50, proves pure codimension three
for the nonregular set in every characteristic.  The finite map
$T\to T/W$ is flat by the same Cohen--Macaulay argument on each
field fibre and the fibrewise flatness criterion over
$\mathbb Z$, so its pullback preserves that codimension.  In
a Borel $B=TU$, no nonregular point lies over the open subset
of $T$ on which every root value differs from one.  Where at
least two independent root equations have value one the base
already has codimension at least two.  At a generic point of
any remaining root divisor the semisimple centralizer has
semisimple rank one.  Successive conjugations in the root
subgroups, dividing only by root values different from one,
reduce the unipotent part to that rank-one centralizer.  It is
nonregular exactly when its one simple-root coordinate is zero,
one further independent equation.  Thus the nonregular set
in $B$ has codimension at least two in every characteristic.
This proves the incidence codimension assertion.  Smoothness
of $\widetilde f$ follows as before from its unipotent-torsor
description.

For a split reductive group with simply connected derived group,
write $C=\operatorname{Rad}(G)$ for its maximal central torus and
use the central isogeny $G_{\rm der}\times C\to G$ with finite
flat central kernel $K$.  This is the radical torus, not the
scheme-theoretic connected centre, which can contain a group such
as $\mu_p$ in characteristic $p$.  The kernel is the graph of
inversion on $G_{\rm der}\cap C$, and its projection to $C$ is
therefore a closed immersion.  In the split case these are
diagonalizable groups, so the restriction
$X^*(C)\to X^*(K)$ is surjective.  In contrast to the
characteristic-zero argument, this cover need not be etale;
it is faithfully flat.  Thus each character of $K$ extends
to a character of the central torus.  Multiplying each
fundamental quotient coordinate by an appropriate torus
character makes it $K$-invariant.  It follows explicitly that
the quotient is an affine space bundle (indeed, after these
choices, a product) over the quotient central torus, and that
the map on quotients is a $K$-torsor.  The group square is
cartesian: the natural equivariant map from its upper left corner
to the pullback is a map of $K$-torsors, hence an isomorphism.
The same argument applies to the universal Borel, since a central
isogeny identifies flag schemes.  Regularity is preserved: the
central isogeny preserves the dimension of each centralizer and
the rank.  Thus the flatness, geometric normality and smoothness
statements descend along this faithfully flat comparison;
the incidence isomorphism descends as well.  Its projectivity
has the independent flag-incidence construction.  This
argument covers inseparable central kernels such as the one
in $\operatorname{SL}_p\times\mathbb G_m\to\operatorname{GL}_p$.

Etale descent of split reductive models gives the theorem for
arbitrary forms.  Every step of the divisor and character-sheaf
proof of D.010 used only normality of the base, these incidence
properties, and the affine-space-times-torus local model; it
did not use characteristic zero.  The same proof therefore
establishes its assertions in this broader scope.
\end{proof}

\subsection{The letter's dimension and subregular questions}
\begin{saresult}{D.018}{Known}{The dimension formula, with equidimensionality}
Let $G$ be a connected complex reductive group of rank $r$, and let
$x$ belong to $G$ or its Lie algebra.  Define
\[
 \mathcal B_x=\{B\in D:x\in B\},\qquad
 \mathcal B_x=\{B\in D:x\in\operatorname{Lie}B\},
\]
respectively, with their reduced variety structures.  Every
irreducible component of $\mathcal B_x$ has dimension
\[
 \dim\mathcal B_x=\tfrac12(\dim Z_G(x)-r).
                                                        \tag{D.18.1}
\]
Thus the dimension conjectured in the letter is $i$ when
$\dim(Gx)=\dim G-r-2i$.  Equidimensionality for group elements
holds over any algebraically closed field.  The cited original
dimension theorems must not, by contrast, be silently restated
with no characteristic hypotheses.
\end{saresult}
\begin{proof}
For group elements, Steinberg's Theorem 4.6, p.217, proves
the equality for unipotent elements under his stated hypothesis
``classical, or characteristic zero or sufficiently large'';
Theorem 6.5, p.223, gives the equality for all elements under
the same hypothesis \cite{D-Steinberg1976}.  In particular the
complex case satisfies it.  His unrestricted inequality is
Theorem 3.5, p.213.  These are the appropriate precise locators;
Proposition 3.1 on p.212 is not the dimension-equality theorem.
Spaltenstein proves equidimensionality in his thesis,
II.1.12, pp.47--48, for unipotent elements, and II.1.17, p.51,
for every element.  Its II.2.5--2.7, pp.53--55, proves the
inequality and equality under explicit hypotheses, including
characteristic zero \cite{D-SpaltensteinThesis}.  Thus this
second source verifies the equality independently in the
scope used here and supplies the stronger assertion about
every component.

Here is the reduction to the Lie-algebra statement, including
its dimension accounting.  Write $x=s+e$ with commuting
semisimple and nilpotent parts, and put $L=Z_G(s)$, a connected
reductive subgroup of rank $r$.  A Borel containing $x$ in its
Lie algebra contains $s$ and $e$ in that Lie algebra: this
follows by the polynomial construction of the Jordan parts
in any faithful representation.  The fixed flag variety of
$s$ is a finite disjoint union of flag varieties of $L$.
Indeed, choose $s\in\mathfrak t$; each Borel through $s$
contains a maximal torus of $L$, and conjugating that torus
inside $L$ to $T$ identifies its component with
$L/(L\cap wBw^{-1})$.  These are flag varieties of $L$.
Intersecting with the condition on $e$ gives a finite disjoint
union of copies of the Springer fibre for $e$ in $L$.
The finite polynomial maps $e\mapsto\exp(e)$ and
$u\mapsto\log(u)$ on nilpotent and unipotent matrices are
inverse, respect every Borel subgroup, and have the same
centralizers.  The group theorem applied to $L$ and $\exp(e)$
therefore gives
\[
 \dim\mathcal B_x=\tfrac12(\dim Z_L(e)-r)
                =\tfrac12(\dim Z_G(x)-r),
\]
and gives equidimensionality.  The same Jordan reduction for
group elements is explicitly recorded by Steinberg in
Lemma 6.2 and Remark 6.3, pp.222--223.
\end{proof}

\begin{saresult}{D.019}{Known}{The subregular differential rank in all simple complex types}
For a simple complex Lie algebra of rank $r$, there is one
subregular nilpotent orbit.  At its points the differential
of $\chi:\mathfrak g\to\mathfrak t/W$ has rank $r-1$ in
\emph{every} simple type.  Uniqueness and this rank do not
characterize the simply laced types.  The corresponding
assertion for the group quotient holds at subregular
unipotent elements of a simply connected complex group.
\end{saresult}
\begin{proof}
Brieskorn states uniqueness of the subregular unipotent class
for all simple groups on p.280, attributing the result to
Tits and Steinberg, and states the rank $r-1$ consequence in
the simply laced case on p.283 \cite{D-Brieskorn}.
The precise surface input in the other types is the
subregular-slice classification: its types are
\[
\begin{array}{c|cccc}
\text{Lie algebra}&B_r&C_r&F_4&G_2\\ \hline
\text{surface}&A_{2r-1}&D_{r+1}&E_6&D_4.
\end{array}
\]
Slodowy states this complete table in the original introduction,
p.IV \cite{D-Slodowy}, for both the nilpotent and unipotent
transverse slices.  That original page has now been opened and
visually checked.  His preceding p.III states the simply laced
identification.  The later attribution to IV.8.3 in
Fu--Juteau--Levy--Sommers, sections 1.4.2 and 1.5.1,
pp.3--4 \cite{D-FJLS}, is consistent with these original statements.
The proof pages of IV.8.3 remain unopened; the present citation
is to the classification actually printed on pp.III--IV.

We give the implication for the rank, rather than treating
it as part of the quoted classification.  The affine
Slodowy slice $e+\mathfrak g^f$ is smooth of dimension
$\dim\mathfrak g^e=r+2$.  Its scheme-theoretic intersection
with $\chi^{-1}(0)$ is one of the displayed surface
singularities, which has tangent-space dimension three at
its singular point: its defining equation in three
variables has no linear term.  Therefore the kernel of
the differential of the restricted quotient has dimension
three, and its rank is $(r+2)-3=r-1$.
The $\mathfrak{sl}_2$ decomposition gives
$\mathfrak g=[\mathfrak g,e]\oplus\mathfrak g^f$;
$d\chi_e$ vanishes on $[\mathfrak g,e]$.  Hence its full
rank is the restricted rank just computed.  In the group
case the complex analytic exponential is an equivariant
local isomorphism at a nilpotent element (its differential
is an invertible polynomial in the nilpotent operator
$\operatorname{ad}e$), and induces a local isomorphism
between the quotient bases near $0$ and $1$.  It transfers
the rank statement.  D.012 supplies a fully explicit
$C_2$ counterexample to the letter's opposite strict
inequality without requiring this classification.
\end{proof}

\begin{saresult}{D.020}{Refuted}{An outer-invariant distinguished class need not exist}
For adjoint type $A_2$ over $\mathbb C$, the nontrivial outer
automorphism interchanges the two classes of Kostant
quadruples.  Thus neither class has a nontrivial full
automorphism, and there is no class invariant under all
outer automorphisms.  In particular no rule functorial
under all group automorphisms can select a single class
in every simple type.
\end{saresult}
\begin{proof}
Use the explicit pair of tori and orbit enumeration in
D.016.  Index the Fourier eigenvectors by $j\in\mathbb Z/3$
and the standard coordinate vectors by $i\in\mathbb Z/3$.
Their standard bilinear pairings identify the dual of
the $j$-th Fourier line with the $-j$-th, whereas every
coordinate line is self-dual.  Thus transpose-inverse
takes a pair of ordered flags $(a,b)\in S_3\times S_3$
to $(\tau a w_0,bw_0)$, where $\tau:j\mapsto-j$ and
$w_0$ reverses the three positions.  Both $\tau$ and
$w_0$ are odd.  The invariant
$\operatorname{sgn}(a)\operatorname{sgn}(b)$ is therefore
negated.  D.016 proved that this sign distinguishes
exactly the two inner orbits.  The outer automorphism
swaps them, as is also checked on all 36 flag pairs by
\texttt{code/rigkos\_a2.py}.  A quadruple's automorphism
group injects into $\operatorname{Out}(A_2)=C_2$ because
its inner stabilizer is trivial.  Neither class is fixed,
so both automorphism groups are trivial.  Any natural
choice of one class would have to be fixed by every
automorphism, which is impossible.
\end{proof}

\begin{saresult}{D.021}{Proved}{Integral cyclic regularity with only $h$ inverted}
For an adjoint simple split group over $R=\mathbb Z[1/h]$,
every vector in
$\bigoplus_{i=0}^r\mathfrak g_{\alpha_i}$ whose $r+1$
coordinates are units is regular semisimple, fibrewise and
etale locally conjugate up to a scalar to one fixed such
vector.  The cyclic-line and Coxeter-normalized definitions
of apposition coincide over $R$ and after any base change.
\end{saresult}
\begin{proof}
We first prove a root-datum assertion over
$R_h=R[\zeta_h]$: a Coxeter transformation on the Cartan
has a primitive-$h$ eigenvector on which every root has
unit value.  In type $A_r$ put $h=r+1$ and use diagonal
coordinates $(1,\zeta,\ldots,\zeta^{h-1})$, whose sum
is zero.  Root values are differences of distinct
powers of $\zeta$.  In types $B_r,C_r$ use
$(1,\zeta,\ldots,\zeta^{r-1})$ for a negative cycle
of length $r$, $h=2r$.  Root values are nonzero sums
or differences of these coordinates, single coordinates,
or twice a coordinate.  In type $D_r$ use
$(1,\zeta,\ldots,\zeta^{r-2},0)$, with $h=2r-2$,
for the negative cycle on the first $r-1$ coordinates
and a negative last coordinate.  All root values are
units: for $0<j<h$, $1-\zeta^j$ is a unit after
inverting $h$, since the product of its Galois
conjugates has only primes dividing $h$.

For the five exceptional types the attached exact
certificate uses the integral cocharacter matrix
$C=c^{-\mathsf t}$ and
\[
 v=\sum_{j=0}^{h-1}\zeta^{-j}C^j e_1.
\]
It verifies $Cv=\zeta v$ modulo the cyclotomic
polynomial and evaluates every root.  The complete
lists, not just their norms, are saved in
\texttt{logs/primitive\_eigenvector\_TYPE.json}.
Exact resultant computation gives the following sets
of absolute norms of root values:
\[
\begin{array}{c|l}
G_2&12,36\\
F_4&144,576\\
E_6&36,144\\
E_7&1728,5184,15552\\
E_8&518400.
\end{array}
\]
Every prime in each row divides the corresponding
$h$.  A cyclotomic integer whose norm is a unit in
$R$ is a unit in $R_h$, by its monic characteristic
equation.  This proves the assertion in all exceptional
types.  The algorithm generates the entire root set
from the displayed Cartan matrices; it tests every
root and retains each polynomial and resultant.
No test in finitely many residue characteristics is
being substituted for this integral certificate.

Now use $n$ and $t=\rho^\vee(\zeta_h)$ from D.015,
whose etale local conjugacy was proved independently
of cyclic regularity.  The vector $v$ is in
$\operatorname{Lie}T_0$, is regular semisimple, and
has primitive eigenvalue under $n$.  Conjugating to
$t$ places it in the eigenspace
\[
 \mathfrak g(1)=\bigoplus_{i=0}^r\mathfrak g_{\alpha_i}.
\]
If any of its coordinates vanished in a geometric
fibre, the remaining roots could all be placed in
a strictly positive half-space, as in D.006.
Its adjoint operator would then be nilpotent.
This contradicts the nonzero root eigenvalues of
$\operatorname{ad}v$.  Its coordinates are therefore
units on the etale conjugating cover.

Finally the action of $T$ on projective cyclic
coordinates is the torus homomorphism with characters
$\alpha_i+\theta$ for $1\le i\le r$.  In the
simple-root basis its determinant is
$1+\sum_i m_i=h$.  It is therefore a finite etale
isogeny over $R$, and acts transitively etale locally
on all choices of nonzero projective coordinates.
Every such cyclic vector is consequently conjugate
up to a unit scalar to the one just constructed,
and is regular semisimple.  Its neutral centralizer
is the corresponding maximal torus, normalized
by $\rho^\vee(\mu_h)$ with Coxeter action.
Conversely the construction from a Coxeter lift
has provided this cyclic vector.  These arguments
are equivariant and all their maps commute with
base change, proving the asserted equivalence of
definitions and the statement over arbitrary $R$-schemes.
\end{proof}

\subsection{Source concordance and the scope of the answers}
The following comparison separates a mathematical answer from an
unverified bibliographical identification.  A theorem named as an
input above has been opened at its displayed locator.
\begin{description}
\item[123/004--010: incidence and quotient.] Brieskorn,
pp.279--280, prints the complex simultaneous-resolution
diagram and credits Grothendieck for its global construction,
Springer for its unipotent restriction.  Steinberg 1976,
Theorem 1.1, pp.209--210, proves the unipotent resolution
when the universal covering is separable; Remark 6.1,
p.222, proves the necessity of that hypothesis for his
construction.  On p.210 he explicitly corrects an earlier
statement for adjoint groups in Springer's \emph{1969}
paper.  That paper must not be confused with the requested
IHES \emph{1966} article.  Slodowy's opened section 4.1,
pp.43--44, repeats the corrected separability hypothesis.
D.009 and D.017 give the relative comparisons and descent;
D.001 explains why the unrestricted Lie-algebra variant fails.
\item[123/011--014: Picard groups.] EGA IV 21.4.13,
Errata et Addenda liste 3, p.361, is the precise generic-fibre
descent result used in the notes.  The number ``53'' is its
erratum number, not a page or a replacement theorem number.
D.010 proves the character-lattice correction, restriction
isomorphisms, and sequences over the normal locally noetherian
base; D.011 proves the projectivity obstruction under its
explicit smallness hypothesis.
\item[Letter: rigidification and opposition.] The finite
quotient and counts are D.004 and D.015.  The splitness
question in the folder is answered negatively by
D.016.  D.013 proves existence of opposite Borels and
refutes the assertion for every pair.  A quadruple has
trivial inner automorphism group; its full automorphism
group is exactly the stabilizer of its class under
$\operatorname{Out}(G)$, via the injective homomorphism
$\operatorname{Aut}(G,T,B,T',B')\to\operatorname{Out}(G)$.
Surjectivity onto this stabilizer follows because an outer
automorphism fixing the class can be uniquely corrected
by an inner conjugation.  This gives a canonical finite
classification through the explicitly constructed finite
quotient, without asserting that its stabilizers are all
the same.  The unqualified phrase ``distinguished in some
sense'' specifies no selection predicate; it cannot be
turned into a uniqueness theorem by choosing a predicate
afterwards.  Its necessary outer-invariance condition is
already impossible in $A_2$, by D.020.
\item[Letter: dimensions and rational singularities.]
D.018 derives the dimension equality and component
statement with the hypotheses of the original sources.
Hesselink's Theorem A, p.249, proves rational resolution
of the characteristic-zero nilpotent variety; D.014 gives
the Jordan-decomposition and descent argument extending
it to all the quotient fibres requested in the letter.
The point is a theorem about $R^i\pi_*\mathcal O$, not a
consequence of dimensions alone.
\item[Letter: subregular uniqueness and rank.] D.019
gives the corrected all-type answer.  Brieskorn's
pp.282--283 give the simply laced universal-deformation
statement.  Slodowy's original introduction, pp.III--IV,
states the complete transverse-slice classification and
explains that in the non-simply-laced cases the deformation
is semiuniversal with fixed diagram symmetries.
His opened section 5.1, pp.60--61, supplies the transversality
and separability conventions for slices.  Appendix I,
pp.152--159, discusses forms over perfect fields of
characteristic zero or sufficiently high characteristic;
it explicitly defers details, so it is not used to extend
our assertions to arbitrary positive characteristics.
\item[Kostant 1959 and 1963.] The 1959 cyclic-element and
Coxeter correspondence is used at its original Theorem 8.6,
p.1026.  In the 1963 \emph{American Journal} paper, distinct
from its shorter \emph{Bulletin} announcement, section 4.7
Theorem 7, p.381, identifies the principal slice with the
adjoint quotient; Theorem 8, p.382, gives its unique
intersection with every regular orbit, and Theorem 9 on
the same page characterizes regularity by independence of
the invariant differentials.  These mathematical passages
were read in the original-paper reprint \cite{D-Kostant1963}.
They do not by themselves tabulate the differential rank
on every nonregular orbit.  The exact source-access ledger
is \path{agents/lie_theory/source_audit_D.md}.
\item[Spaltenstein LNM 946.] The public preview of
Chapter II.1, pp.42--43, was opened: it announces
connectedness and equidimensionality and begins the
projective-line argument.  The subsequent proof pages
were not available through that preview.  The independently
opened original 1977 thesis, at the locators in D.018,
is the primary input actually used for the dimension proof.
\item[Springer IHES 30, 1966.] Section 5.7 distinguishes
regular from smoothly regular nilpotents.  Theorem 5.9,
p.138, states the good-characteristic and root-system
fundamental-group condition for smoothly regular
nilpotents.  Section 4.12, pp.134--135, shows that a
regular unipotent centralizer need not be connected in
bad characteristic.  These are reasons to retain the
isogeny and prime hypotheses; this article is not a
proof of the arbitrary-base claims in the folder.
\item[Steinberg LNM 366.] In addition to the definitions
on pp.76--77, the original appendix, pp.156--158, is now
opened.  It relates the Kleinian resolutions to the
subregular group slices and states the corresponding
deformation theorem.  Its displayed standalone hypotheses
omit an ADE restriction, although the preceding ADE
discussion may indicate the intended scope.  D.027 proves
that the unrestricted reading fails already for
$\mathrm{Sp}_4$.  Slodowy p.IV gives the required
fixed-symmetry formulation outside ADE.
The body of III.3.10 remains unopened and is not used as
an inspected theorem.  The comparisons just made depend
on the actual appendix passages, not on that inaccessible
section.
\end{description}

\subsection{The full characteristic-zero Jacobian question in the letter}

\begin{saresult}{D.022}{Proved}{All nilpotent quotient ranks, including the exceptional types}
Let $\mathfrak g$ be a simple complex Lie algebra, and let
$\chi:\mathfrak g\longrightarrow\mathfrak t/W$ be its adjoint quotient.
The following formulas determine $\operatorname{rank}(d\chi_e)$ for every
nilpotent $e$.  In the classical types let
$\lambda=(\lambda_1\ge\cdots\ge\lambda_k)$ be its Jordan partition in
the defining representation, and put $m=\lambda_1$.
\[
\begin{array}{c|c}
\text{type and condition}&\operatorname{rank}(d\chi_e)\\\hline
 A_r&m-1\\
 B_r,\ C_r&\lfloor m/2\rfloor\\
 D_r,\ k>2&\lfloor m/2\rfloor\\
 D_r,\ \lambda=(2r-i,i),\ i\text{ odd}&(2r-i+1)/2\\
 D_r,\ r\text{ even},\ \lambda=(r,r)&r/2.
\end{array}                                                     \tag{D22.1}
\]
The orthogonal and symplectic partition parity conditions are understood;
the two very-even $D_r$ orbits have the same rank.
For the exceptional types, the following table lists every orbit whose
rank is at least two.  The rank is the number of displayed exponents.
Every other nonzero nilpotent orbit has rank one; the zero orbit has
rank zero.  Complete weighted diagrams, representatives, and exact
certificates for all $5+16+21+45+70=157$ orbits, including zero,
are supplied in the accompanying machine-readable files.
\end{saresult}

\begin{longtable}{@{}lp{29mm}p{91mm}@{}}
\toprule Type&Nonzero exponents&Nilpotent orbit labels\\\midrule\endhead
$G_{2}$&$1,\allowbreak 5$&$G_{2}$\\
$F_{4}$&$1,\allowbreak 5$&$C_{3},\allowbreak\;B_{3},\allowbreak\;F_{4}(a_{2})$\\
$F_{4}$&$1,\allowbreak 5,\allowbreak 7$&$F_{4}(a_{1})$\\
$F_{4}$&$1,\allowbreak 5,\allowbreak 7,\allowbreak 11$&$F_{4}$\\
$E_{6}$&$1,\allowbreak 4$&$A_{4},\allowbreak\;A_{4}+A_{1}$\\
$E_{6}$&$1,\allowbreak 5$&$D_{4}$\\
$E_{6}$&$1,\allowbreak 4,\allowbreak 5$&$A_{5},\allowbreak\;D_{5}(a_{1}),\allowbreak\;E_{6}(a_{3})$\\
$E_{6}$&$1,\allowbreak 4,\allowbreak 5,\allowbreak 7$&$D_{5}$\\
$E_{6}$&$1,\allowbreak 4,\allowbreak 5,\allowbreak 7,\allowbreak 8$&$E_{6}(a_{1})$\\
$E_{6}$&$1,\allowbreak 4,\allowbreak 5,\allowbreak 7,\allowbreak 8,\allowbreak 11$&$E_{6}$\\
$E_{7}$&$1,\allowbreak 5$&$D_{4},\allowbreak\;(A_{5})^{\prime\prime},\allowbreak\;D_{4}+A_{1},\allowbreak\;D_{5}(a_{1}),\allowbreak\;(A_{5})^{\prime}$\\
$E_{7}$&$1,\allowbreak 5$&$A_{5}+A_{1},\allowbreak\;D_{5}(a_{1})+A_{1},\allowbreak\;D_{6}(a_{2}),\allowbreak\;E_{6}(a_{3}),\allowbreak\;E_{7}(a_{5})$\\
$E_{7}$&$1,\allowbreak 5$&$A_{6}$\\
$E_{7}$&$1,\allowbreak 5,\allowbreak 7$&$D_{5},\allowbreak\;D_{5}+A_{1},\allowbreak\;D_{6}(a_{1}),\allowbreak\;E_{7}(a_{4}),\allowbreak\;E_{6}(a_{1})$\\
$E_{7}$&$1,\allowbreak 5,\allowbreak 7,\allowbreak 9$&$D_{6},\allowbreak\;E_{7}(a_{3})$\\
$E_{7}$&$1,\allowbreak 5,\allowbreak 7,\allowbreak 11$&$E_{6}$\\
$E_{7}$&$1,\allowbreak 5,\allowbreak 7,\allowbreak 9,\allowbreak 11$&$E_{7}(a_{2})$\\
$E_{7}$&$1,\allowbreak 5,\allowbreak 7,\allowbreak 9,\allowbreak 11,\allowbreak 13$&$E_{7}(a_{1})$\\
$E_{7}$&$1,\allowbreak 5,\allowbreak 7,\allowbreak 9,\allowbreak 11,\allowbreak 13,\allowbreak 17$&$E_{7}$\\
$E_{8}$&$1,\allowbreak 7$&$D_{5},\allowbreak\;D_{5}+A_{1},\allowbreak\;D_{6}(a_{1}),\allowbreak\;E_{7}(a_{4}),\allowbreak\;E_{6}(a_{1})$\\
$E_{8}$&$1,\allowbreak 7$&$D_{5}+A_{2},\allowbreak\;D_{6},\allowbreak\;D_{7}(a_{2}),\allowbreak\;A_{7},\allowbreak\;E_{6}(a_{1})+A_{1}$\\
$E_{8}$&$1,\allowbreak 7$&$E_{7}(a_{3}),\allowbreak\;E_{8}(b_{6}),\allowbreak\;D_{7}(a_{1}),\allowbreak\;E_{8}(a_{6})$\\
$E_{8}$&$1,\allowbreak 7,\allowbreak 11$&$E_{6},\allowbreak\;E_{6}+A_{1},\allowbreak\;E_{7}(a_{2}),\allowbreak\;D_{7},\allowbreak\;E_{8}(b_{5})$\\
$E_{8}$&$1,\allowbreak 7,\allowbreak 11$&$E_{8}(a_{5})$\\
$E_{8}$&$1,\allowbreak 7,\allowbreak 11,\allowbreak 13$&$E_{7}(a_{1}),\allowbreak\;E_{8}(b_{4}),\allowbreak\;E_{8}(a_{4})$\\
$E_{8}$&$1,\allowbreak 7,\allowbreak 11,\allowbreak 13,\allowbreak 17$&$E_{7},\allowbreak\;E_{8}(a_{3})$\\
$E_{8}$&$1,\allowbreak 7,\allowbreak 11,\allowbreak 13,\allowbreak 17,\allowbreak 19$&$E_{8}(a_{2})$\\
$E_{8}$&$1,\allowbreak 7,\allowbreak 11,\allowbreak 13,\allowbreak 17,\allowbreak 19,\allowbreak 23$&$E_{8}(a_{1})$\\
$E_{8}$&$1,\allowbreak 7,\allowbreak 11,\allowbreak 13,\allowbreak 17,\allowbreak 19,\allowbreak 23,\allowbreak 29$&$E_{8}$\\
\bottomrule
\end{longtable}


\begin{proof}
\par\medskip\noindent\textbf{Classical types.}\quad
Use the nondegenerate trace pairing to identify differentials with
matrices.  In type $A_r$, basic invariants may be taken to be
$\operatorname{tr}(X^j)$ for $2\le j\le r+1$.  At a nilpotent $e$
their gradients are the nonzero scalar multiples $j e^{j-1}$.
The matrices $e,e^2,\ldots,e^{m-1}$ are linearly independent: apply a
putative relation to a cyclic vector in a block of length $m$.
Higher powers vanish.  This proves the first row.

For $B_r,C_r$ take the even power traces, of degrees $2,4,\ldots,2r$.
Their gradients are the odd powers of $e$, which remain in the
orthogonal or symplectic Lie algebra.  The trace form restricted to
that Lie algebra is nondegenerate in characteristic zero.  Exactly
$\lfloor m/2\rfloor$ odd powers are nonzero, and the same cyclic-vector
argument makes them independent.  This proves the next row.

For $D_r$ use the even power traces of degrees $2,4,\ldots,2r-2$
and the Pfaffian.  Its differential is zero exactly when the
skew-symmetric matrix has rank at most $2r-4$: its partial derivatives
are, up to sign, the Pfaffians of the complementary $(2r-2)$-minors.
At rank $2r-2$ its gradient is a nonzero bivector on the two-dimensional
kernel.  Thus the Pfaffian contributes nothing if $k>2$.
If $k=2$ and both blocks have odd lengths, the defining orthogonal
space is an orthogonal direct sum of the two nondegenerate cyclic
blocks.  All odd powers of $e$ are block diagonal, whereas the
Pfaffian gradient is the nonzero cross term between their kernel
lines.  It is independent of those powers and contributes one.
If both block lengths are even, the partition condition forces
$\lambda=(r,r)$ with $r$ even.  The two cyclic spaces are totally
isotropic and paired nondegenerately.  Both the Pfaffian gradient
and $e^{r-1}$ are nonzero skew endomorphisms with image in $\ker e$.
The space of such endomorphisms identifies with
$\Lambda^2\ker e$ and is one-dimensional.  They are proportional,
so the Pfaffian contributes no extra dimension.  This proves all
rows of (D22.1), including the repeated exponent in even type $D$.
These formulas agree with Richardson's Theorems 4.2.1 and 4.3.3,
p.233, and Theorem 4.4.2, p.234 \cite{D-Richardson1987}.

\par\medskip\noindent\textbf{Exact construction in the exceptional types.}\quad
Here are sufficient details to make the finite computation a proof
with inspectable inputs, rather than a sample of orbit ranks.
For the simply laced root lattice, let $A$ be its symmetric Cartan
matrix, and put
\[
 \epsilon(a,b)=(-1)^{\sum_i a_i b_i+
                      \sum_{i>j}a_i b_j A_{ij}}.
\]
In the basis $\{E_a:a\in\Phi\}\cup\{h_1,\ldots,h_r\}$ use
\[
\begin{aligned}
 [h_i,E_a]&=(Aa)_i E_a,& [h_i,h_j]&=0,\\
 [E_a,E_b]&=\epsilon(a,b)E_{a+b}&& (a+b\in\Phi),\\
 [E_a,E_{-a}]&=\epsilon(a,-a)\sum_i a_i h_i,
\end{aligned}                                                     \tag{D22.2}
\]
with all other brackets zero.  The cocycle is bimultiplicative;
its commutator is $(-1)^{(a,b)}$.  Formula (D22.2) is skew, and
the code checks Jacobi on \emph{every} ordered-independent basis
triple.  There are respectively $76076$, $383306$, and $2511496$
triples for $E_6,E_7,E_8$.  The Cartan action and nonzero opposite-root
brackets identify the resulting split simple Lie algebras with these
root systems.  Thus no unverified sign table enters the computation.

The $F_4$ algebra is the fixed algebra of the pinned $E_6$ involution,
whose simple-node permutation in Bourbaki numbering is
$(1\ 6)(3\ 5)$.  Its ordered simple-node orbits are
$\{2\},\{4\},\{3,5\},\{1,6\}$.
The $G_2$ algebra is the fixed algebra of pinned triality on $D_4$,
with ordered orbits $\{2\},\{1,3,4\}$.
The signs on nonsimple root vectors are obtained recursively from
(D22.2), not assumed to be positive.  Taking invariant sums in each
root orbit gives dimensions $52$ and $14$; the resulting Cartan
matrices, all brackets, and all $22100$ and $364$ Jacobi triples
are checked.  In our $G_2$ order the first simple root is long;
this reverses the order in the input Atlas/SLA diagrams.

For $E_6$ construct the $27$-dimensional minuscule module by the
Weyl orbit of $\omega_1$.  In the weight basis, a simple positive
root operator has coefficient one when it raises a weight, and
zero otherwise; the corresponding negative root operator has
coefficient minus one in the cocycle convention (D22.2).
Cartan operators are diagonal.  Recursion through (D22.2) determines
all root operators.  The program checks the representation identity
$\rho([a,b])=[\rho(a),\rho(b)]$ on all $3003$ distinct basis pairs.
This verifies the representation rather than merely relying on
its stated dimension.

For each nonzero weighted diagram $D=(D_i)$, compute
$h=\sum_i c_i h_i$ from $A^{\mathsf t}c=D$.
Choose an integral $e\in\mathfrak g(2)$ and solve the rational
linear system $[e,f]=h$ with $f\in\mathfrak g(-2)$.
The saved certificate verifies all three $\mathfrak{sl}_2$ relations.
Consequently it proves that the chosen $e$ has precisely the claimed
weighted diagram.  Random integer choices are used only to find
such an $e$; success is certified by exact equalities, and failure
is never interpreted as nonexistence.  The orbit dimension is
independently checked from
\[
 \dim\mathfrak g^e=\dim\mathfrak g(0)+\dim\mathfrak g(1),
\]
which follows by decomposing the adjoint module into irreducible
$\mathfrak{sl}_2$ modules.
The complete lists of nonzero characteristics are cross-checked
against the author's SLA lists, with the indicated root-order
change for $G_2$ and before the software permutation for $F_4$.
Completeness and the classification by weighted diagrams are the
standard classification stated in de Graaf, \S2, pp.281--282;
his Appendix A gives the complete exceptional lists
\cite{D-deGraaf}. The label and dimension input is the complete
author-hosted Adams table \cite{D-AdamsOrbitTables}, reopened and
checked byte-for-byte against the saved input. No assertion of completeness is inferred from
successful random searches.

Let $\rho$ be the adjoint representation in $G_2,F_4,E_7,E_8$
and the just-verified minuscule representation in $E_6$.
For each fundamental degree $d$, use the invariant
$F_d(X)=\operatorname{tr}(\rho(X)^d)$.  We calculate, for every
basis vector $Y$,
\[
 (dF_d)_e(Y)=d\operatorname{tr}
                 \bigl(\rho(e)^{d-1}\rho(Y)\bigr).       \tag{D22.3}
\]
Only vectors of $h$-weight $-2(d-1)$ can contribute; this vanishing
follows directly by conjugating by the $\mathfrak{sl}_2$ torus.
The remaining traces are exact integers.  Distinct nonzero
$dF_d$ have distinct $h$-weights, so they are independent.

There is also a check that these particular traces generate the
\emph{entire} invariant algebra.  The fundamental degrees are distinct
in every exceptional type.  In a homogeneous basic generating set,
$F_d$ is a scalar multiple of the degree-$d$ generator plus a
polynomial in lower-degree generators.  At a nilpotent element,
the differential of every product of positive-degree invariants
vanishes.  At the regular nilpotent representative our computation
finds every $dF_d$ nonzero.  Hence each of those scalar multiples
is nonzero; induction on degree proves that the $F_d$ themselves
are basic generators.  Thus (D22.3) computes the quotient rank
exactly, including its upper bound.

The complete run uses
\path{agents/lie_theory/code/exceptional_jacobian.py},
\path{agents/lie_theory/code/all_exceptional_jacobians.py}, and
\path{agents/lie_theory/code/exceptional_orbits.json}.
Its full output is
\path{agents/lie_theory/logs/all_exceptional_jacobians.txt};
the JSON companion saves every triple and every nonzero derivative.
The timed rerun completed all $157$ orbits in $22.197$ seconds.
The displayed table was generated from that complete output, not
transcribed from a secondary table.

\par\medskip\noindent\textbf{The formerly undecided $E_8(a_2)$ entry.}\quad
Richardson's paper explicitly leaves exponent $19$ for $E_8(a_2)$
undecided on p.235 and in its Appendix, p.239.
Our representative has diagram $22202022$ and maximum $h$-weight
$38$.  The certificate gives
\[
 (dF_{20})_e(E_{-\theta})=3617744486400\ne0.              \tag{D22.4}
\]
The nonzero exponents are exactly $1,7,11,13,17,19$, so the rank
is six.  Exponents $23$ and $29$ vanish already by the weight bound.
A separate implementation using SymPy sparse exponentiation by
squaring verifies (D22.4), in $3.293$ seconds:
\path{agents/lie_theory/code/verify_jacobian_independent.py} and
\path{agents/lie_theory/logs/jacobian_independent.txt}.
This is an independent arithmetic check, not a claim that a new
historical result is being established here.

The original scan of Richardson p.239 contains two further literal
errors: its regular $E_8$ row ends in $27$ instead of $29$, and
its $E_8(a_1)$ row includes a spurious $9$.
These are printed-source errors, not OCR substitutions.  They
contradict the fundamental degrees and the subregular rank, and
are corrected by the exact computations above.
\end{proof}

\begin{saresult}{D.023}{Proved}{Reduction of every complex Lie-algebra rank to D.022}
Let $\mathfrak g$ be a complex reductive Lie algebra, and write
$x=s+e$ for its commuting Jordan decomposition.  Put
$\mathfrak l=\mathfrak g^s$ and decompose
$[\mathfrak l,\mathfrak l]=\bigoplus_j\mathfrak l_j$ and
$e=\sum_j e_j$.  For the typed quotient maps
$\chi_{\mathfrak g}:\mathfrak g\to\mathfrak t/W$ and
$\chi_j:\mathfrak l_j\to\mathfrak t_j/W_j$, one has
\[
 \operatorname{rank}(d\chi_{\mathfrak g})_x
   =\dim Z(\mathfrak l)+
      \sum_j\operatorname{rank}(d\chi_j)_{e_j}.            \tag{D23.1}
\]
For a simply connected complex semisimple group $G$, the same
formula determines the rank of $\chi_G:G\to T/W$ at every
$g=su$: replace $\mathfrak l$ by $\operatorname{Lie}Z_G(s)$
and $e$ by $\log u$.  Formula (D23.1) is not, without a further
finite quotient, a formula for arbitrary isogeny types.
\end{saresult}
\begin{proof}
In the Lie-algebra case $L=Z_G(s)$ is connected, and its Weyl group
is the stabilizer $W_s$.  On the complement of $\mathfrak l$,
$\operatorname{ad}s+\operatorname{ad}e$ is invertible: the first
summand is invertible semisimple, and the second is a commuting
nilpotent.  The conjugation map from the associated slice
$G\times^L(s+\mathfrak l)$ to $\mathfrak g$ is therefore etale
at $[1,s+e]$.  The induced map on quotients is etale at the point
under consideration: locally the finite Weyl quotient near $s$
is the quotient by its stabilizer $W_s=W_L$.
This can also be seen on completed local rings by taking the
product of completions at the points of the finite $W$-orbit
and then taking invariants; characteristic zero makes that
operation exact.  On $\mathfrak l$ the quotient splits as the
identity on its center and the product of its simple quotients.
This proves (D23.1).  It is also Richardson's Proposition 1.2,
p.229 \cite{D-Richardson1987}, with all central dimensions shown.

For the group statement, $Z_G(s)$ is connected when $G$ is simply
connected.  The original source explicitly states this in
Steinberg, Remark 2.10, p.53 \cite{D-Steinberg1965}.
Here is a root-theoretic justification of the needed fact.  Write
$s=\exp(u+iv)$ in a torus, with $u,v$ in its real coroot space,
and move $u$ into a closed Weyl chamber.  Since the cocharacter
lattice is the coroot lattice, a Weyl element fixing $s$ fixes $u$
and carries $v$ to $v$ modulo $2\pi$ times the coroot lattice.
The subsystem perpendicular to $u$ is generated by a subset of
the simple roots; its coroot lattice is the intersection of its
real span with the full coroot lattice.  Hence the relevant
stabilizer is a point stabilizer in that subsystem's affine
Weyl group.  Such a stabilizer is generated by the affine
reflections through the point: the incident alcoves form the
finite chamber system of those walls, and reflection across
a separating wall reduces the number of separating walls.
Their linear parts are exactly the reflections with
$\alpha(s)=1$.  Thus $W_s=W(\Phi_s)$.  Bruhat decomposition
then shows that the centralizer is generated by its torus and
these root groups, so it is connected (compare Steinberg 2.8,
p.53).  The same slice argument applies, and exponential
is an equivariant analytic local isomorphism from a neighbourhood
of $\log u$ to one of $u$ inside that centralizer.  Its differential
is invertible because $\operatorname{ad}(\log u)$ is nilpotent.
Exponential identifies the corresponding local torus quotients.
The analytic differential rank is the algebraic differential rank,
so (D23.1) applies without imposing a false global isomorphism
between the group and its Lie algebra.
\end{proof}

\begin{saresult}{D.024}{Proved}{The finite-center correction for every complex group isogeny}
Let $q:\widetilde G\to G=\widetilde G/Z$ be a central isogeny of
complex semisimple groups, with $\widetilde G$ simply connected.
Let $F_i=\operatorname{tr}(\rho_i)$ be its fundamental characters,
$1\le i\le r$, and let $\omega_i|_Z$ be their central characters.
For a lift $\widetilde g$ of $g$, set $b=(F_i(\widetilde g))$ and
\[
 \Gamma_b=\{z\in Z:\omega_i(z)F_i(\widetilde g)
                         =F_i(\widetilde g)\text{ for all }i\},
 \quad I_b=\{i:\omega_i|_{\Gamma_b}=1\}.
\]
For any basis $Y_1,\ldots,Y_{\dim G}$ of
$\widetilde{\mathfrak g}=\operatorname{Lie}\widetilde G$, the exact rank
at $g$ is
\[
 \operatorname{rank}(d\chi_G)_g
  =\operatorname{rank}\left(
     \operatorname{tr}\bigl(\rho_i(\widetilde g)
                                  d\rho_i(Y_j)\bigr)
                    \right)_{i\in I_b,\ 1\le j\le\dim G}.     \tag{D24.1}
\]
This is a finite matrix formula for every element and every isogeny
in every simple type.  Fundamental representations and their central
characters are determined by the root datum.  If $\Gamma_b=1$ it
reduces to D.023; in particular this holds for unipotent $g$.
\end{saresult}
\begin{proof}
The fundamental characters identify
$\widetilde G/\!/\widetilde G$ with $\mathbb A^r$
(Steinberg, Theorem 6.1(b), p.62; the finite-center quotient
is also stated in 6.18, p.66 \cite{D-Steinberg1965}).
Translation by $z\in Z$ acts diagonally on these coordinates as
$F_i\mapsto\omega_i(z)F_i$.  Taking invariant coordinate rings gives
\[
 G/\!/G\simeq(\widetilde G/\!/\widetilde G)/Z.
\]
The quotient locally at $b$ is the quotient by its stabilizer
$\Gamma_b$.  In translated coordinates $u_i=F_i-b_i$, this
stabilizer still acts diagonally.  The linear terms of its invariant
power series are exactly the span of those $u_i$ with $i\in I_b$.
Therefore the differential of this finite quotient retains exactly
these coordinate functionals.  Higher-degree invariant generators
have zero differential at the origin, even when the quotient is
singular.  As $q$ is etale in characteristic zero, it does not
change the rank.  Differentiating $F_i$ in the right-translated
tangent direction $Y_j$ gives precisely the trace in (D24.1).
This proves the formula and its independence of the chosen lift
and tangent basis.

For a unipotent lift, $F_i(\widetilde g)=\dim\rho_i\ne0$.
The fundamental weights generate the character group of the
center, so their common kernel on $Z$ is trivial.  Thus
$\Gamma_b=1$ in this case.  This explains exactly why the
nilpotent/unipotent comparison is isogeny-independent in
characteristic zero, while the all-element statement is not.

For the smallest discrepancy, let
$\widetilde G=\mathrm{SL}_2$, $G=\mathrm{PGL}_2$, and
$\widetilde g=\operatorname{diag}(i,-i)$.
The simply connected quotient coordinate is $t=\operatorname{tr}$;
the adjoint quotient coordinate is $t^2$.  At this lift $t=0$,
$dt\ne0$, but $d(t^2)=0$.  Thus the ranks are respectively
one and zero at corresponding regular semisimple elements.
The exact symbolic check, with its output and timing, is
\path{agents/lie_theory/code/group_quotient_rank.py} and
\path{agents/lie_theory/logs/group_quotient_rank.txt}.
\end{proof}

\begin{saresult}{D.025}{Proved}{Every quadruple automorphism group in type A}
Let $G=\mathrm{PGL}_n(\mathbb C)$, $n\ge3$, and let
$Q=(T,B,T',B')$ be an ordered apposition quadruple.
Its automorphism group is either $1$ or $C_2$.
If $n$ is odd it is always trivial; in particular no inner-conjugacy
class of quadruples is fixed by the outer involution.
If $n=2k$ is even, put $m=n/2^{v_2(n)}$ and let
$\omega(m)$ be the number of its distinct prime factors
(with $\omega(1)=0$).  Among the
\[
 C_n=\frac{(n!)^2}{n^2\varphi(n)}
\]
inner-conjugacy classes, precisely
\[
 F_n=\frac{(2^k k!)^2\,2^{\omega(m)}}{n\varphi(n)}
       \begin{cases}1,&4\mid n,\\[2pt]1/2,&n\equiv2\pmod4\end{cases}
                                                               \tag{D25.1}
\]
have automorphism group $C_2$; the other $C_n-F_n$ have trivial
automorphism group.  The constructive test in the proof determines
the stabilizer of each class, not only its aggregate number.
For $A_3$ the two numbers are $8$ and $10$, out of $18$ classes.
For $A_1$ the outer group is trivial and all stabilizers are trivial.
\end{saresult}
\begin{proof}
Use the cyclic matrix model of D.016 in dimension $n$.
Let $T$ be the diagonal torus with coordinate basis indexed by
$\mathbb Z/n$, and let $T'$ be its Fourier transform, with Fourier
columns $(\zeta^{ij})_i$, also indexed by $j\in\mathbb Z/n$.
A Borel through either torus is a total ordering of its eigenlines,
so the choices are pairs $(p,q)\in S_n\times S_n$ of ordered lists.
The common inner normalizer acts by
\[
 (p,q)\longmapsto(ap+x,a^{-1}q+y),\quad
 a\in(\mathbb Z/n)^\times,\ x,y\in\mathbb Z/n.           \tag{D25.2}
\]
Indeed cyclic permutation matrices translate $p$, diagonal powers
of $(1,\zeta,\ldots,\zeta^{n-1})$ translate $q$, and the permutation
$i\mapsto ai$ acts by $a$ and $a^{-1}$ on the two lists.
These give $n^2\varphi(n)$ distinct common-normalizer elements;
D.003 proves that this is the full order.  The action on pairs
of orderings is free.  This also proves the displayed formula
for $C_n$ directly in this model.

Transpose-inverse induces the outer involution
\[
 \tau(p,q)=(p^{\rm rev},-q^{\rm rev}).                   \tag{D25.3}
\]
For the first list this is duality of a complete coordinate flag.
For the second use $F^{-\mathsf t}=n^{-1}FJ$, where $F$ is the
Fourier matrix and $J$ negates its column index.  Hence (D25.3)
is the actual outer action, including its sign.
A class has nontrivial automorphism precisely when there exist
$a,x,y$ in (D25.2) such that
\[
 p^{\rm rev}=ap+x,\qquad q^{\rm rev}=-a^{-1}q-y.         \tag{D25.4}
\]
If this occurs, the lift of the outer element fixing the quadruple
is unique, because its inner stabilizer is trivial.  Its square
is therefore the identity.  This gives the promised constructive
test and the abstract stabilizer $C_2$.

Let $R$ denote reversal of an ordered list.  Each of the two
affine permutations in (D25.4) must be conjugate to $R$.
If $n$ is odd, it must have exactly one fixed point.  Since the
first affine permutation is an involution, $a^2=1$.
Modulo each odd prime power, $a$ is either $1$ or $-1$.
Exactly one fixed point forces $a-1$ to be invertible, hence
$a=-1$ modulo every prime power and therefore modulo $n$.
The second affine permutation is then a translation; an
involutive translation of an odd cyclic group is the identity,
which has $n>1$ fixed points.  This contradiction proves the
odd-$n$ assertion uniformly, strengthening D.020.

Suppose now $n=2k$.  Reversal has no fixed points, and its
centralizer in $S_n$ has order $c_n=2^k k!$.
Let $S_n^{\rm aff}$ be the number of triples $(a,x,y)$ for which
both affine maps in (D25.4) are fixed-point-free involutions.
For each such triple there are exactly $c_n^2$ pairs $(p,q)$
satisfying (D25.4).  Conversely, each outer-fixed inner class
contains $n^2\varphi(n)$ pairs, each with exactly one such triple.
Consequently
\[
 F_n=\frac{S_n^{\rm aff}c_n^2}{n^2\varphi(n)}.           \tag{D25.5}
\]
We compute $S_n^{\rm aff}$ by the Chinese remainder theorem.
At an odd prime power $\ell^b\Vert m$, the slope has two choices
$\pm1$.  If it is $1$, the involution condition forces $x=0$
and permits all $\ell^b$ choices of $y$; if it is $-1$ the roles
are reversed.  Both affine involutions always have fixed points
on this odd-primary factor.  Multiplying these contributions gives
$m2^{\omega(m)}$.

On $\mathbb Z/2^a$, fixed-point-freeness of the first map means
\[
 (a_0+1)x=0,\qquad x\notin(a_0-1)(\mathbb Z/2^a),
\]
where $a_0$ denotes its slope (to distinguish it from the exponent
$a$); the second condition is obtained by replacing the slope by
$-a_0$.  For $a=1$ there is one slope and one admissible pair of
translations.  For $a\ge2$, the slopes $a_0=1,-1$ each yield
$2^{a-1}$ pairs: the translation of the slope-one map is the
unique nonzero element of order two, and that of the slope-minus-one
map is odd.  For $a\ge3$ the only other roots of $a_0^2=1$ are
$1+2^{a-1}$ and $-1+2^{a-1}$; in either case the first displayed
condition already forces $x$ into the forbidden image, so they
contribute zero.  A product of affine involutions is fixed-point-free
exactly when at least one primary factor has no fixed point;
here only the $2$-primary factor can do so.  It follows that
\[
 S_n^{\rm aff}=
 \begin{cases}
 n2^{\omega(m)},&4\mid n,\\
 (n/2)2^{\omega(m)},&n\equiv2\pmod4.
 \end{cases}
\]
Substitution into (D25.5) proves (D25.1).

The finite checks are exhaustive over all pairs for $A_2,A_3,A_4$,
and independently count the affine involutions just used:
\begin{verbatim}
A2: 36 pairs; 18 normalizer elements; 2 classes; C2:0; trivial:2.
A3: 576 pairs; 32 normalizer elements; 18 classes; C2:8; trivial:10.
A4: 14400 pairs; 100 normalizer elements; 144 classes; C2:0; trivial:144.
\end{verbatim}
The program and full representatives are
\path{agents/lie_theory/code/type_a_outer_stabilizers.py} and
\path{agents/lie_theory/logs/type_a_outer_stabilizers.json};
the timing and output are in the corresponding text log.
These finite checks supplement the uniform proof, which handles
all ranks rather than extrapolating from these three cases.
\end{proof}

\begin{saresult}{D.026}{Proved}{All remaining quadruple automorphism groups}
For ordered apposition quadruples in a simple adjoint complex group,
the following list, together with D.025, determines every possible
automorphism group and the number of inner-conjugacy classes of each type.
\begin{itemize}
\item In $B_r,C_r,G_2,F_4,E_7,E_8$, every automorphism group is trivial.
\item In $D_r$, $r\ge5$, put
$N_r=|W(D_r)|^2/(8(r-1)\varphi(2r-2))$.
Exactly $2N_r/r^2$ classes have automorphism group $C_2$;
the other $N_r-2N_r/r^2$ have trivial automorphism group.
\item In $D_4$, among the $768$ classes, $504$ have trivial
automorphism group, $252$ have group $C_2$, and $12$ have group $S_3$.
None has group $C_3$.  The quotient by all automorphisms has $180$ classes.
\item In $E_6$, among the $18662400$ classes, $27648$ have group
$C_2$ and $18634752$ have trivial automorphism group.
The quotient by all automorphisms has $9345024$ classes.
\end{itemize}
The finite certificates also determine the stabilizer of an individual
class by the exact twisted-Weyl criterion in the proof.
\end{saresult}
\begin{proof}
\medskip\noindent\textbf{Equivariant identification of the two Weyl torsors.}
Let $\Gamma$ be the pinned diagram-automorphism group.
In types $D_r,E_6$, the cyclic element
$x=\sum_i x_{\alpha_i}+x_{-\theta}$ and the principal element
$t=\rho^\vee(\zeta)$ can be chosen fixed by $\Gamma$; the sign
calculation D.005 verifies this on the highest-root vector.
Thus $T$ and $T'=C_G(x)$ are $\Gamma$-stable and $\Gamma$ fixes
the regular semisimple vector $x\in\operatorname{Lie}T'$.
The fixed subspace of the real cocharacter space of $T'$ contains
a regular vector: otherwise it would be contained in some root
hyperplane, and its complexification could not contain $x$.
More explicitly, finitely many proper real root hyperplanes
cannot cover that real vector space.  Choose its chamber;
the associated Borel $B'_0$ is $\Gamma$-stable.
Together with the pinned Borel $B_0\supset T$, this identifies
the two Borel sets with Weyl groups on which $\Gamma$ acts by
diagram automorphisms.  The action on either diagram has the
same outer class as its action on $G$, and is faithful.

The Coxeter element $c$ induced by $t$ on $T'$ commutes with
$\Gamma$.  We may identify it with the following standard
$\Gamma$-invariant Coxeter element.  In $D_r$ with the end-node
involution, use a negative $(r-1)$-cycle and a negative singleton
in signed coordinates.  Every invariant element in the $D_r$
Coxeter conjugacy class is conjugate to this one by the fixed
Weyl subgroup: commuting with the sign change on the final
coordinate forces the unique singleton of the signed cycle type
to occupy that coordinate.  The remaining negative cycle is
unique up to signed permutation of the other coordinates.
For full triality in $D_4$, take
$c=s_1s_3s_4s_2$ in Bourbaki numbering.  Complete enumeration
finds precisely two triality-fixed elements of its Weyl conjugacy
class, and they are a single orbit under $W^{S_3}$.
For $E_6$, take the bipartite product
$c=(s_1s_4s_6)(s_2s_3s_5)$.
Complete enumeration finds precisely $96$ diagram-fixed elements
of its Weyl conjugacy class, and they are a single orbit under
$W^{C_2}$; these cardinalities are certified, not assumed from
a folding picture.  This establishes the equivariant comparison
needed to use the following finite model for the geometric action.

Write $A=N_W(\langle c\rangle)$ and
$B_t=\{b\in W:b(t)=t^a,\ a\in U_h\}$.
By D.003 the common normalizer is
$\Phi=B_t\times_{U_h}A$, acting freely on $W\times W$ by
left multiplication in the two factors.  The outer group acts
by $(w,w')\mapsto(\gamma(w),\gamma(w'))$.
The stabilizer of the class of $(w,w')$ consists exactly of
those $\gamma\in\Gamma$ for which
\[
 \bigl(\gamma(w)w^{-1},\gamma(w')w'^{-1}\bigr)\in\Phi.
                                                               \tag{D26.1}
\]
As the inner stabilizer is trivial, this subgroup is precisely
the automorphism group of the corresponding quadruple.
This also supplies the individual-class test.

For an involution $\sigma$, define
$\Delta_\sigma=\{\sigma(w)w^{-1}:w\in W\}$.
Every fibre of $w\mapsto\sigma(w)w^{-1}$ has cardinality
$|W^\sigma|$: equality for $w_1,w_2$ is equivalent to
$w_2^{-1}w_1\in W^\sigma$.
Counting (D26.1) before dividing by the free $\Phi$-action gives
\[
 \#\operatorname{Fix}(\sigma;\Phi\backslash(W\times W))
 =\frac{|W^\sigma|^2}{|\Phi|}
   \#\bigl(\Phi\cap(\Delta_\sigma\times\Delta_\sigma)\bigr).
                                                               \tag{D26.2}
\]
This counts every point of the finite quotient, even when that
quotient is much larger than $W$.

\medskip\noindent\textbf{The uniform $D_r$ calculation.}
Put $m=r-1\ge3$, so $h=2m$ and $z=4$.
Use signed-coordinate matrices.  The end-node involution is
conjugation by the sign change in the last coordinate.
Its coboundary set consists exactly of the identity and the
$r-1$ diagonal matrices changing signs at coordinate $j$ and
the last coordinate.  Also
$|W^\sigma|=2^{r-1}(r-1)!$.
Let $\rho^\vee=(m,m-1,\ldots,1,0)$.
A nontrivial such diagonal matrix belongs to $B_t$ only if
it changes the coordinate with value $m$; its exponent is then one.
Indeed the zero last coordinate excludes the half-integral
coset of $hP^\vee$, and an unchanged coordinate of value one,
or unchanged coordinates of values two and three, force
$a=1\pmod {2m}$.  The changed value must then be a multiple of $m$.
Thus $B_t\cap\Delta_\sigma$ has two elements, both of exponent one.

On the other hand $A\cap\Delta_\sigma=\{1\}$.
Conjugating the negative $m$-cycle by one coordinate sign change
alters precisely two of its edge signs.  Its underlying permutation
is unchanged, so equality with $c^a$ requires $a\equiv1\pmod m$.
The only alternative to $a=1$ is $a=m+1$, available only for even
$m$; it changes all $m$ edge signs.  Since $m\ge3$, neither
alternative equals the change of exactly two edge signs.
Consequently the intersection in (D26.2) has two elements.
Formula (D26.2) gives
$2|W^\sigma|^2/|\Phi|=2N_r/r^2$.
For $r\ge5$ the outer group is $C_2$, which proves the asserted
full classification.  The exact finite checks for $4\le r\le20$
are in \path{agents/lie_theory/code/type_d_outer_stabilizers.py}, with complete output
in \path{agents/lie_theory/logs/type_d_outer_stabilizers.txt}; the uniform argument,
not extrapolation, proves the statement for every rank.

\medskip\noindent\textbf{Triality in $D_4$.}
The program enumerates all $192$ Weyl matrices, all $12$ elements
of $A$, all $8$ elements of $B_t$, and all $48$ elements of $\Phi$.
It then enumerates all $768$ free orbits in $W\times W$ and
checks all six diagram automorphisms on each orbit.
The stabilizer orders occur with multiplicities
\[
\begin{array}{c|rrrr}
\text{order}&1&2&3&6\\\hline
\text{number of classes}&504&252&0&12.
\end{array}
\]
As an independent count, for each $\gamma\in S_3$ the program
sums over $(b,a)\in\Phi$ the product of the solution counts
$\gamma(w)=bw$ and $\gamma(w')=aw'$.
After division by $48$, the fixed-class counts are
$768$ for the identity, $96$ for each transposition, and $12$
for each three-cycle.  These agree with the complete class list.
They give $180$ full-automorphism orbits by averaging.
The matrices, class representatives, and actual stabilizer subgroups
are all in \path{agents/lie_theory/logs/d4_outer_stabilizers.json}; the two methods
and the Coxeter-class comparison run in
\path{agents/lie_theory/code/d4_outer_stabilizers.py}.
The full text output is
\path{agents/lie_theory/logs/d4_outer_stabilizers.txt};
the recorded run takes $0.227$ seconds.

\medskip\noindent\textbf{The $E_6$ calculation.}
The program enumerates all $51840$ Weyl matrices.
It finds $|W^\sigma|=1152$ and exactly $45$ coboundaries,
with every fibre counted independently to have size $1152$.
The sets $A,B_t,\Phi$ have cardinalities $48,12,144$.
Exactly three pairs in $\Phi$ have both entries in
$\Delta_\sigma$; their second entries are all the identity.
Thus (D26.2) gives
\[
 \frac{3\cdot1152^2}{144}=27648.
\]
Subtracting from $51840^2/144=18662400$ gives the stated trivial
stabilizer count.  Every tested matrix, every coboundary with a
witness, the three contributing pairs, and the fixed-Coxeter
conjugacy check are saved in \path{agents/lie_theory/logs/e6_outer_stabilizers.json}.
The program is \path{agents/lie_theory/code/e6_outer_stabilizers.py},
with text output \path{agents/lie_theory/logs/e6_outer_stabilizers.txt}.
The complete run, including the geometric-comparison certificate,
takes $4.801$ seconds.  The other simple types have trivial outer
group, so D.004 already makes their quadruple automorphism groups
trivial.  This completes all simple complex types.
\end{proof}

\begin{saresult}{D.027}{Refuted}{The unrestricted reading of the printed deformation statement}
The hypotheses printed immediately before Steinberg's
LNM~366 Appendix theorem, pp.157--158, are insufficient
for ordinary versality in part~(a): the simple simply
connected group $\mathrm{Sp}_4(\mathbb C)$ is a
counterexample.  This is a correction to the standalone
wording.  The preceding discussion of ADE singularities
may indicate an intended restriction to those types.
\end{saresult}
\begin{proof}
The original pp.157--158 were opened and checked against
their page images.  They specify a simple simply connected
complex group $G$, a subregular unipotent element $u$,
and a transverse slice of dimension $r+2$, then assert
that the restriction of
$p:G\longrightarrow T/W$ is a universal deformation
of the corresponding Kleinian singularity
\cite[Appendix, pp.157--158]{D-Steinberg1974}.
We test the weaker necessary condition of versality, so
the historical use of \emph{universal} for a semiuniversal
deformation does not affect the counterexample.

Take $G=\mathrm{Sp}_4(\mathbb C)$ with symplectic matrix
\[
 J=\begin{pmatrix}0&I_2\\-I_2&0\end{pmatrix},\qquad
 N=\begin{pmatrix}0&I_2\\0&0\end{pmatrix},\qquad
 u=I_4+N .
\]
The equations $N^{\mathsf t}J+JN=0$ and
$N^{\mathsf t}JN=0$ prove $u\in G$, and $N^2=0$
proves that $u$ is unipotent.  A general Lie-algebra
element has the form
$X=\left(\begin{smallmatrix}A&B\\C&-A^{\mathsf t}\end{smallmatrix}\right)$,
where $B,C$ are symmetric, and
\[
 [X,N]=
 \begin{pmatrix}-C&A+A^{\mathsf t}\\0&C\end{pmatrix}.
\]
Thus its centralizer has dimension $1+3=4=r+2$:
$A$ is skew-symmetric, $C=0$, and $B$ is arbitrary
symmetric.  Centralizers are smooth in characteristic
zero, so this verifies that $u$ is subregular.
The quotient $T/W$ is smooth of dimension two.  Its subregular
transverse surface has type $D_3=A_3$, by the original
table in Slodowy's introduction, p.IV
\cite{D-Slodowy}; thus its completed local ring is
\[
 R=\mathbb C[[x,y,z]]/(F),\qquad
 F=x^4+y^2+z^2 .
\]
Here the formal completion is enough: versality in
analytic germs implies versality for first-order
deformations.  Put $A=\mathbb C[\epsilon]/(\epsilon^2)$.
In a flat first-order deformation over $A$ of
$R$, lift the three coordinate generators.  The resulting
map from $A[[x,y,z]]$ onto the deformation ring is
surjective by the complete Nakayama lemma.  Flatness
identifies the reduction of its kernel with $(F)$;
lifting that generator and applying Nakayama again shows
that the kernel has one generator $F+\epsilon g$.
Conversely, in the quotient by $F+\epsilon g$, an
equation $\epsilon q=(F+\epsilon g)h$ first forces
$h=\epsilon h_1$, since $F$ is a nonzerodivisor modulo
$\epsilon$.  It then gives
$q=Fh_1+\epsilon q_1=\epsilon(q_1-gh_1)$ in the
quotient.  Therefore the kernel of multiplication by
$\epsilon$ is its image, which is precisely flatness
over $A$.  Thus these are the first-order deformations.
Two are isomorphic by an isomorphism inducing the
identity on the special fibre exactly when their
coefficients differ by
\[
 (F,F_x,F_y,F_z).
\]
Indeed, a lift of the images of $x,y,z$ under any such
isomorphism gives substitutions
$x_i\mapsto x_i+\epsilon a_i$ in the power-series ring,
and a change of defining equation multiplies by
$1+\epsilon b$.  Expanding shows that the difference
is $\sum a_iF_{x_i}+bF$.  Conversely these substitutions
and this multiplier realize every such difference.
They are invertible because they are the identity
modulo $\epsilon$.
Consequently the tangent space of deformation classes is
\[
 T^1_R=
 \mathbb C[[x,y,z]]/(F,4x^3,2y,2z)
 =\mathbb C[[x]]/(x^3),
 \qquad \dim_{\mathbb C}T^1_R=3 .
\]
Every family over the smooth two-dimensional germ
$(T/W,p(u))$ has a linear first-order class map
\[
 T_{p(u)}(T/W)\longrightarrow T^1_R
\]
with image of dimension at most two: its pullbacks to
$\operatorname{Spec}A$ are obtained by substituting the
two first-order base parameters.  Versality would make
this map surjective, which is impossible.

For a reproducible exact finite check,
\path{agents/lie_theory/code/check_steinberg_universality.py}
computes the reduced rational Groebner basis
$(x^3,y,z)$ and the three standard monomials
$1,x,x^2$.  The actual output is saved in
\path{agents/lie_theory/logs/check_steinberg_universality.txt}:
\begin{quote}\ttfamily\small
TJURINA DIMENSION: 3\\
ADJOINT-QUOTIENT BASE TANGENT DIMENSION: 2\\
PASS: the unrestricted C2 family cannot be versal or universal\\
RUNTIME\_SECONDS: 0.070802045
\end{quote}
The mathematical correction is supplied by the same
original Slodowy p.IV: in types $B,C,F_4,G_2$ the
semiuniversal property concerns deformations with the
specified diagram symmetries.  No proof from the
unopened sections IV.8.3 or III.3.10 is attributed to
either original book.
\end{proof}
